\specialchapter{引言}

谱理论旨在研究某些赋范代数的性质，其基本的例子是：拓扑空间上复值连续函数所成的代数、Banach 空间的自同态所成的代数，以及关于局部紧群上 Haar 测度可积的函数所成的卷积代数。仅由这一简单的说明未必能看出这项研究的范围，但要认识其重要性，只需注意：两位当代数学家若不属于同一学派，则除了线性代数以及微分学与积分学的初等部分之外，他们共同具有的知识往往只是 Fourier 分析的某些方面。

在本书的头两章中，我们展开谱理论中最基本的思想，并更详细地考察上面提到的第三个例子，即局部紧交换群框架下的 Fourier 变换。

随后的各章将首先继续紧线性映射的一般理论，然后是 Hilbert 空间上算子的理论，后一情形还包括仅定义在这类空间的一个子空间（通常是稠密子空间）上的线性映射。众所周知，这一理论是量子力学数学形式体系的基础。最后，我们将给出拓扑群的酉表示的初等理论，特别是紧群的酉表示；在《李群与李代数》一书第九章中我们已有机会用到它。

头两章的标题与本书 1967 年出版的第一版相同。它们经过了修订、更正与扩充。

为便于阅读，我们在下面指出本版中最重要的改动，特别是增补的内容。

关于第 I 章：

- 在 §3 的第 3 小节中，我们引入真偏映射的概念，它使局部紧空间与交换 C*-代数之间的 Gelfand 对应能够以函子化的方式表述。

- 在 §4 中，我们增加了第 13 和第 14 小节，它们讨论实或复可赋范代数中的全纯函数演算，以及无单位元的代数中的全纯函数演算。

- 在 §6 中，连续函数演算的叙述已重新改写（第 6 小节）。第 9 至第 11 小节是新增加的；它们讨论一般 C*-代数中正元素的概念、此类代数的近似单位元，以及 C*-代数对闭双边理想的商。

- §7 论及 Banach 空间自同态的最基本的谱性质，它基本上是新写的；第 7、8 两小节部分地取材于原 §6 的某些小节。正是这一新的一节将作为第 IV 章中更深入地研究 Hilbert 空间自同态的基础。

关于第 II 章：

- 关于 Fourier 变换的 §1 已作了细致的修订，例如讨论 Pontryagin 对偶的函子性的第 6 小节。在新的第 9 小节中，我们把对群 \(\mathbf{R}^n\) 与 \((\mathbf{R}/\mathbf{Z})^n\) 而言最重要的那些命题陈述得更为明确。\(\mathbf{R}^n\) 上的 Fourier 理论，其最自然的框架在于缓增分布理论，这里仅通过其本质性质加以陈述，而不给出证明；这些证明将在第 IV 章中给出。

- 我们最后增加了一张 Fourier 理论一览表，其中汇集了这一理论所涉及的主要函数空间的定义和性质，以及基本的公式。

习题中有许多是新增加的。我们特别指出第 II 章 §1 中习题数量的显著增加，这反映了 Fourier 理论的数学重要性。

\chapter{赋范代数}

除非另有说明，本章所考虑的代数都假定是结合的。

\section{谱与特征标}

在本节中，字母 K 表示一个交换域。若 E 与 F 是 K-向量空间，我们记 \(E \otimes F = E \otimes_{K} F\) 。

\subsection{幺代数}

K 上的幺代数是指一个序偶 \((A,e)\)，其中 A 是一个带单位元的 K 上的代数，而 e 是 A 的单位元。由于 e 由 A 唯一确定，我们有时会滥用语言，径直说 A 是一个幺代数。若 \((A,e)\) 与 \((A',e')\) 是幺代数，则 \((A,e)\) 到 \((A',e')\) 的一个保单位元态射是指一个态射 \(\varphi\)，它把 A 映到 \(A'\)，且使得 \(\varphi(e)=e'\)。\((A,e)\) 的一个幺子代数是指一个序偶 \((A',e)\)，其中 \(A'\) 是 A 的包含 e 的子代数。

我们常将单位元记作 1。

引理 1. --- 设 A 是一个代数。对 A 的每个幂等元 j，A 的子空间 jAj 是使得 xj = jx = x 的 \(x \in A\) 所成的集合。它是 A 的一个子代数，以 j 为单位元。

证明是初等的。

\subsection{幺代数中元素的谱}

定义 1. --- 设 A 是 K 上的一个幺代数，e 为其单位元。对每个 \(x \in \mathrm{A}\)，我们把 \(x\) 相对于 A 的谱称为由 \(\lambda \in \mathrm{K}\) 中使 \(\lambda e - x\) 不可逆的那些元素所成的集合。

x 的谱将记作 \(\mathrm{Sp}_{\mathrm{A}}(x)\)，若不致混淆，也记作 \(\mathrm{Sp}(x)\)。\(\mathrm{Sp}_{\mathrm{A}}(x)\) 在 K 中的补集称为 x 的预解集。

注. --- 1) 若 A = \{0\}，则有 Sp(0) = ∅。

\begin{enumerate}
\def\labelenumi{\arabic{enumi})}
\setcounter{enumi}{1}
\item
  若 \(A \neq \{0\}\)，则我们有 \(\mathrm{Sp}(\lambda e) = \{\lambda\}\)，其中 \(\lambda \in K\) 任意。
\item
  要使 \(x \in \mathrm{A}\) 可逆，必须且只需 \(0 \notin \operatorname{Sp}(x)\)。
\item
  设 \(R \in K(X)\) 是一个有理分式，并设 \(x \in A\) 是一个可在 R 中代入的元素，也就是说 (A, IV, p.~20)，存在 P 与 \(Q \in K[X]\)，使得 \(R = P/Q\) 且 \(Q(x)\) 可逆；于是可以构成 A 的元素 \(R(x) = P(x) \cdot Q(x)^{-1} = Q(x)^{-1} \cdot P(x)\)；它不依赖于 P 与 Q 的选取。我们有 \(0 \notin Q(\mathrm{Sp}(x))\)，从而 \(\mathrm{Sp}(x)\) 的每个元素都可在 R 中代入。
\end{enumerate}

我们有 \(\mathrm{R}(\mathrm{Sp}(x)) \subset \mathrm{Sp}(\mathrm{R}(x))\)。事实上，设 \(\lambda \in \mathrm{Sp}(x)\)；存在多项式 \(P_{1}\)，使得 \(\mathrm{R}(\lambda) - \mathrm{R}(\mathrm{X}) = (\lambda - \mathrm{X})\mathrm{P}_{1}(\mathrm{X}) / \mathrm{Q}(\mathrm{X})\)；于是 \(\mathrm{R}(\lambda)e - \mathrm{R}(x) = (\lambda e - x)(\mathrm{P}_{1}(x)/\mathrm{Q}(x))\)，因而 \(\mathrm{R}(\lambda) - \mathrm{R}(x)\) 不可逆，故 \(\mathrm{R}(\lambda) \in \mathrm{Sp}(\mathrm{R}(x))\)。

反之，设域 K 是代数闭的。先设 R 不是常数，我们来证明 \(\mathrm{Sp}(\mathrm{R}(x)) = \mathrm{R}(\mathrm{Sp}(x))\)。设 \(\mu \in \mathrm{Sp}(\mathrm{R}(x))\)。由于 R 不是常数，P - \(\mu Q\) 不是零多项式；设 \(\mu Q - P = \alpha \prod (\lambda_i - X)\) 是一个一次因式的分解，于是 \(\mu e - R(x) = \alpha \prod (\lambda_i e - x) Q(x)^{-1}\)。由于 \(\mu e - R(x)\) 不可逆，存在 i 使 \(\lambda_i e - x\) 不可逆，因而 \(\lambda_i \in \mathrm{Sp}(x)\)，于是 \(\mathrm{R}(\lambda_i) = \mu \in \mathrm{R}(\mathrm{Sp}(x))\)。

当 R 为常数时，等式 \(\mathrm{Sp}(\mathrm{R}(x)) = \mathrm{R}(\mathrm{Sp}(x))\) 同样成立，只要 \(\mathrm{Sp}(x)\) 非空。

\begin{enumerate}
\def\labelenumi{\arabic{enumi})}
\setcounter{enumi}{4}
\item
  设代数 A 非零。设 \(x \in A\) 是一个幂零元。设 n 是使 \(x^{n} = 0\) 的一个整数。\(x^{n}\) 的谱化为 0，从而由注 4，x 的谱亦然。
\item
  设 A 与 B 是 K 上的幺代数，\(\varphi\colon\mathrm{A}\to\mathrm{B}\) 是一个保单位元态射。对每个 \(x\in\mathrm{A}\)，有 \(\mathrm{Sp}_{\mathrm{B}}(\varphi(x))\subset\mathrm{Sp}_{\mathrm{A}}(x)\)。
\item
  设 A 是一个幺代数，R 是其根 (A, VIII, p.~150, 定义 2)，\(\varphi\) 是 A 到 \(B = A/R\) 上的典范态射。若 \(x \in A\)，则有 \(\mathrm{Sp}_{\mathrm{B}}(\varphi(x)) = \mathrm{Sp}_{\mathrm{A}}(x)\)。事实上，只需证明：若 \(\varphi(x)\) 在 B 中可逆，则 \(x\) 在 A 中可逆。而今，若 \(y \in A\) 使得 \(\varphi(x)\varphi(y) = \varphi(y)\varphi(x) = \varphi(e)\)，则有 \(xy \in e + R\)，\(yx \in e + R\)，因而 \(xy\) 与 \(yx\) 均可逆 (A, VIII, p.~151, 定理 1)，从而 \(x\) 可逆。特别地，若 \(x \in R\)，则我们有 \(\mathrm{Sp}_{\mathrm{A}}(x) = \{0\}\)，只要 \(A \neq \{0\}\)。
\item
  设 \((\mathrm{B}_{i})_{i\in\mathrm{I}}\) 是一族幺代数，其中 \(\mathrm{B}_{i}=(\mathrm{A}_{i},e_{i})\)（\(i\in I\)）。令 \(A=\prod_{i}A_{i}, e=(e_{i})_{i\in\mathrm{I}}\)。则 \((\mathrm{A},e)\) 是一个幺代数，称为诸 \(B_{i}\) 的积。若 \(x=(x_{i})_{i\in\mathrm{I}}\in\mathrm{A}\)，则有 \(\mathrm{Sp}_{\mathrm{A}}(x)=\bigcup_{i}\mathrm{Sp}_{\mathrm{A}_{i}}(x_{i})\)。
\end{enumerate}

例. --- 1) 设 X 是一个集合，A = K \(^{X}\) 是定义在 X 上而取值于 K 的函数所成的代数。A 的元素 f 的谱就是 f 的值集。

\begin{enumerate}
\def\labelenumi{\arabic{enumi})}
\setcounter{enumi}{1}
\tightlist
\item
  设 A 是 K 上有限秩的幺代数。要使 \(x \in A\) 可逆，必须且只需从 A 到 A 的线性映射 \(y \mapsto xy\) 的行列式非零。由此推出，x 的谱是 x 的特征多项式的根集 (A, III, p.~110)。因此，若 A 是 K 上有限维向量空间 V 的自同态代数，则 x 的谱就是 x 的特征值集。当 V 为无限维时，情况并非总是如此（参见 I, p.~153, 习题 2）。
\end{enumerate}

\subsection{预解式}

定义 2. --- 设 A 是 K 上的一个幺代数，且 \(x \in \mathbf{A}\)。对每个 \(\lambda \in \mathrm{K} - \mathrm{Sp}(x)\)，令

\[
\mathrm{R} (x, \lambda) = (\lambda e - x) ^ {- 1}.
\]

由 \(K-\mathrm{Sp}(x)\) 到 A 的映射 \(\lambda\mapsto\mathrm{R}(x,\lambda)\) 称为 \(x\) 的预解式。

固定 \(x\) 后，\(R(x, \lambda)\) 的各个值两两可交换。若 \(\lambda, \mu \in K\)，我们有：

\[
(\lambda - \mu) e = (\lambda e - x) - (\mu e - x)
\]

因而，若 \(\lambda, \mu \in \mathrm{K} - \mathrm{Sp}(x)\)，则有关系式

\[
(\lambda - \mu) \mathrm{R} (x, \lambda) \mathrm{R} (x, \mu) = \mathrm{R} (x, \mu) - \mathrm{R} (x, \lambda).\tag{1}
\]

若 \(x,y\in \mathrm{A}\) 且 \(\lambda \in \mathrm{K}\)，我们有：

\[
y - x = (\lambda e - x) - (\lambda e - y)
\]

因而，若 \(\lambda\in\mathrm{K}-(\mathrm{Sp}(x)\cup\mathrm{Sp}(y))\)，则有关系式

\[
\mathrm{R} (y, \lambda) (y - x) \mathrm{R} (x, \lambda) = \mathrm{R} (y, \lambda) - \mathrm{R} (x, \lambda).\tag{2}
\]

\subsection{代数中元素的谱}

设 A 是 K 上的一个代数。回顾 (A, III, p.~4)，我们在向量空间 \(\widetilde{\mathbf{A}} = \mathbf{K} \times \mathbf{A}\) 上定义一个代数结构，使得：

\[
(\lambda , a) (\mu , b) = (\lambda \mu , \lambda b + \mu a + a b).
\]

设 \(e = (1, 0)\)。则 \((\tilde{\mathrm{A}}, e)\) 是由 A 添加单位元而得到的幺代数。代数 A 等同于双边理想 \(\{0\} \times A\)，它是 \(\tilde{A}\) 的理想；而且代数 A 交换当且仅当 \(\tilde{A}\) 交换。

若 A' 是 K 上的另一个代数，设 \((\tilde{\mathrm{A}}', e')\) 是由 A' 添加单位元而得到的幺代数，并设 \(\varphi\) 是从 A 到 A' 的一个态射；则存在唯一的一个从 \((\tilde{\mathrm{A}}, e)\) 到 \((\tilde{\mathrm{A}}', e')\) 中的保单位元态射，它延拓 \(\varphi\)。

设 A 是 K 上的一个代数，且 \(x \in A\)。我们把 x 相对于 A 的谱称为 x 相对于由 A 添加单位元而得到的幺代数 \(\widetilde{A}\) 的谱。这个集合将记作 \(\mathrm{Sp}_{\mathrm{A}}^{\prime}(x)\)，或在不致混淆时记作 \(\mathrm{Sp}^{\prime}(x)\)。我们有 \(0 \in \mathrm{Sp}_{\mathrm{A}}^{\prime}(x)\)，它对一切 \(x \in A\) 成立。

若 \(\varphi\) 是从 A 到某个代数 B 中的一个态射，则我们有 \(\mathrm{Sp}_{\mathrm{B}}^{\prime}(\varphi(x)) \subset \mathrm{Sp}_{\mathrm{A}}^{\prime}(x)\)。

注. --- 1) 设 (A,1) 是一个幺代数。若 \(x \in A\)，则有

\[
\operatorname{Sp} _ {\mathrm{A}} ^ {\prime} (x) = \operatorname{Sp} _ {\mathrm{A}} (x) \cup \{0 \}.
\]

事实上可以验证 \((e-1) \cdot A = A \cdot (e-1) = 0\)，于是 \(\widetilde{A}\) 是 A 与 \(K(e-1)\) 的幺代数积。因此我们的断言由 I, p.~3 的注 8 得出。

\begin{enumerate}
\def\labelenumi{\arabic{enumi})}
\setcounter{enumi}{1}
\tightlist
\item
  由注 1 推出，若 B 是 K 上的一个代数且 \(x \in B\)，则有：
\[
\mathrm{Sp} _ {\mathrm{B}} ^ {\prime} (x) = \mathrm{Sp} _ {\widetilde {\mathrm{B}}} (x) = \mathrm{Sp} _ {\widetilde {\mathrm{B}}} (x) \cup \{0 \} = \mathrm{Sp} _ {\widetilde {\mathrm{B}}} ^ {\prime} (x).
\]

\item
  若 x 属于 A 的根 (A, VIII, p.~430, 定义 3)，则 \(\mathrm{Sp}_{\mathrm{A}}^{\prime}(x)=\{0\}\)。这由 I, p.~3 的注 7 得出。
\end{enumerate}

命题 1. --- 设 A 是一个代数，且 \(x, y \in \mathrm{A}\)。我们有 \(\mathrm{Sp}'(xy) = \mathrm{Sp}'(yx)\)。

过渡到 \(\widetilde{\mathbf{A}}\)，我们化归到 A 具有单位元 \(e\) 的情形。此时只需证明：若 \(\lambda \neq 0\) 使得 \(xy - \lambda e\) 有一个逆 \(u\)，则 \(yx - \lambda e\) 可逆。令 \(z = yux - e\)。由于 \(xyu = \lambda u + e\)，有

\[
\begin{array}{r l} (y x - \lambda e) z & = y (x y u) x - y x - \lambda y u x + \lambda e \\ & = y (\lambda u + e) x - y x - \lambda y u x + \lambda e = \lambda e \end{array}
\]

同样地，有 \(z(yx - \lambda e) = \lambda e\)。由于 \(\lambda \neq 0\)，我们看出 \(yx - \lambda e\) 可逆。

若 A 是一个幺代数且 \(x, y \in A\)，则前面的命题蕴涵 \(\operatorname{Sp}(xy) \cup \{0\} = \operatorname{Sp}(yx) \cup \{0\}\)，但可能有 \(\operatorname{Sp}(xy) \neq \operatorname{Sp}(yx)\)（参见 I, p.~153, 习题 3）。

\subsection{全子代数}

设 A 是 K 上的一个幺代数。

定义 3. --- A 的全子代数是指 A 的一个幺子代数 B，使得 B 中每个在 A 内可逆的元素在 B 内也可逆。

换言之，B 是 A 的全子代数，当且仅当 \(\mathrm{Sp}_{\mathrm{B}}(x)=\mathrm{Sp}_{\mathrm{A}}(x)\) 对一切 \(x\in B\) 成立。

A 的一族全子代数的交是 A 的一个全子代数。

设 M 是 A 的一个子集。A 的包含 M 的诸全子代数之交，是 A 的包含 M 的最小全子代数。

它称为由 M 生成的 A 的全子代数。M 在 A 中的交换子 M′ 是 A 的一个全子代数（因为，若 x 在 A 中可逆且与 M 交换，则 \(x^{-1}\) 与 M 交换）。因此，M 的双交换子 M'' 是 A 的一个全子代数，它包含由 M 生成的 A 的全子代数。

若 M 的元素两两可交换，我们有 \(M \subset M'\) 与 \(M'' \subset M'''\)；于是代数 \(M''\) 是交换的，从而由 M 生成的全子代数也是交换的。

A 的一个极大交换子代数是全子代数，因为它等于其交换子。

引理 2. --- 设 \((x_{\lambda})_{\lambda \in \Lambda}\) 是 A 中两两交换的一族元素。由 \((x_{\lambda})\) 生成的全子代数是由形如 \(\mathrm{R}((x_{\lambda}))\) 的元素所成的集合，其中 \(\mathrm{R} \in \mathrm{K}((\mathrm{X}_{\lambda}))\) 取遍使族 \((x_{\lambda})\) 可代入的那些有理分式。

设 B 是由族 \((x_{\lambda})\) 生成的 A 的全子代数，并用 \(B_{1}\) 表示由形如 \(\mathrm{R}((x_{\lambda}))\) 的元素所成的集合，其中 \(\mathrm{R}((\mathrm{X}_{\lambda}))\) 是使 \((x_{\lambda})\) 可代入的有理分式。明确地说，\(B_{1}\) 是 A 中形如 \(\mathrm{P}((x_{\lambda}))\mathrm{Q}((x_{\lambda}))^{-1}\) 的元素所成的集合，其中 \(\mathrm{P},\mathrm{Q}\in\mathrm{K}[(\mathrm{X}_{\lambda})]\) 且 \(\mathrm{Q}((x_{\lambda}))\) 在 A 中可逆。集合 \(B_{1}\) 是 A 的包含族 \((x_{\lambda})\) 的一个幺子代数。它是一个全子代数：若 \(\mathrm{P}((x_{\lambda}))\mathrm{Q}((x_{\lambda}))^{-1}\) 在 A 中可逆，则 \(\mathrm{P}((x_{\lambda}))\) 在 A 中可逆，且逆 \(\mathrm{Q}((x_{\lambda}))\mathrm{P}((x_{\lambda}))^{-1}\)（即 \(\mathrm{P}((x_{\lambda}))\mathrm{Q}((x_{\lambda}))^{-1}\) 的逆元）属于 \(B_{1}\)。于是我们有 \(B\subset B_{1}\)。另一方面，若 \(P,Q\in K[(X_{\lambda})]\)，且 \(\mathrm{Q}((x_{\lambda}))\) 在 A 中可逆，则 \(\mathrm{P}((x_{\lambda}))\in\mathrm{B}\) 且 \(\mathrm{Q}((x_{\lambda}))\in\mathrm{B}\)，因而 \(\mathrm{Q}((x_{\lambda}))^{-1}\in\mathrm{B}\) 且 \(\mathrm{P}((x_{\lambda}))\mathrm{Q}((x_{\lambda}))^{-1}\in\mathrm{B}\)，从而 \(B_{1}\subset B\)。

\subsection{交换幺代数的特征标}

定义 4. --- 设 A 是 K 上的一个交换幺代数。幺特征标是指从 A 到 K 的一个保单位元态射。

在不致混淆时，可简称特征标以代替幺特征标。A 的幺特征标所成的集合记作 X(A)。若 A 是零代数，则 X(A) 为空集。

设 A 与 B 是 K 上的交换幺代数，h 是从 A 到 B 的一个保单位元态射。映射 \(\chi \mapsto \chi \circ h\)（它把 \(\mathsf{X}(\mathsf{B})\) 映到 \(\mathsf{X}(\mathsf{A})\) 中）记作 \(\mathsf{X}(h)\)。若 k 是从 B 到某个交换幺代数的态射，则有 \(\mathsf{X}(k \circ h) = \mathsf{X}(h) \circ \mathsf{X}(k)\)。映射 \(\mathsf{X}(\mathrm{Id}_{\mathrm{A}})\) 是 \(\mathsf{X}(\mathrm{A})\) 上的恒同映射。

若 h 是满射，则 \(\mathsf{X}(h)\) 是从 \(\mathsf{X}(\mathsf{B})\) 到在 h 的核上为零的 A 的特征标所成集合上的一个双射。

设 \((\mathrm{A}_{1},e_{1}),\ldots,(\mathrm{A}_{n},e_{n})\) 是 K 上的交换幺代数，A 是幺代数 \(A_{1}\times\cdots\times A_{n}\)，其单位元为 \((e_{1},\ldots,e_{n})\)。对每个 i，把 \(A_{i}\) 等同于 A 的一个理想，并设 \(\pi_{i}\) 是从 A 到 \(A_{i}\) 上的典范映射。则 \(\mathsf{X}(\pi_{i})\) 是从 \(\mathsf{X}(\mathrm{A}_{i})\) 到集合 \(X_{i}\)（即 A 的在 \(\prod_{j\neq i}A_{j}\) 上为零的特征标所成的集合）上的一个双射。诸集合 \(X_{i}\) 两两不相交。另一方面，设 \(\chi\in\mathsf{X}(\mathrm{A})\)。由于 \(1=\sum\chi(e_{i})\)，存在 i 使 \(\chi(e_{i})\neq0\)。对每个 \(j\neq i\) 及每个 \(y\in A_{j}\)，有 \(\chi(e_{i})\chi(y)=\chi(e_{i}y)=\chi(0)=0\)，因而 \(\chi(\mathrm{A}_{j})=0\)。于是 \(\chi\) 在 \(\prod_{j\neq i}A_{j}\) 上为零，从而 \(\mathsf{X}(\mathrm{A})\) 是诸 \(X_{i}\) 的并。

设 B 是幺代数 \(A_{1} \otimes \cdots \otimes A_{n}\)。用 \(h_{i}\) 表示典范态射 \(A_{i} \to B\)。则

\[
\chi \mapsto (\chi \circ h _ {1}, \dots , \chi \circ h _ {n})
\]

是从 \(\mathsf{X}(\mathsf{B})\) 到 \(\mathsf{X}(\mathsf{A}_1)\times \dots \times \mathsf{X}(\mathsf{A}_n)\) 中的一个映射，而

\[
\left(\chi_ {1}, \dots , \chi_ {n}\right) \mapsto \chi_ {1} \otimes \dots \otimes \chi_ {n}
\]

是从 \(\mathsf{X}(\mathsf{A}_{1}) \times \cdots \times \mathsf{X}(\mathsf{A}_{n})\) 到 \(\mathsf{X}(\mathsf{B})\) 中的一个映射。可以验证这两个映射互为逆双射，据此我们把 \(\mathsf{X}(\mathsf{B})\) 与 \(\mathsf{X}(\mathsf{A}_{1}) \times \cdots \times \mathsf{X}(\mathsf{A}_{n})\) 视为同一。

设 A 是 K 上的一个交换幺代数。设 Y 是 A 的余维数为 1 的理想所成的集合。对每个 \(\chi \in \mathsf{X}(\mathsf{A})\)，有 \(\operatorname{Ker}(\chi) \in \mathsf{Y}\)。映射 \(\chi \mapsto \operatorname{Ker}(\chi)\) 是从 \(\mathsf{X}(\mathsf{A})\) 到 Y 上的一个双射。事实上，若 \(I \in Y\)，则存在唯一一个从幺 K-代数 A/I 到 K 上的同构，而复合态射

\[
\mathrm{A} \longrightarrow \mathrm{A} / \mathrm{I} \longrightarrow \mathrm{K}
\]

是 A 的以 I 为核的唯一特征标。

定义 5. --- 设 A 是 K 上的一个交换幺代数。对每个 \(x \in \mathrm{A}\)，我们用 \(\mathcal{G}_{\mathrm{A}}(x)\)（或简作 \(\mathcal{G}(x)\)）表示映射 \(\chi \mapsto \chi(x)\)（从 \(\mathrm{X}(\mathrm{A})\) 到 K）。它称为 \(x\) 的 Gelfand 变换。

映射 \(\mathcal{G}\) 是从 A 到幺代数 \(\mathrm{K}^{\mathsf{X}(A)}\)（由 \(\mathsf{X}(A)\) 到 K 的映射所成的幺代数）中的一个保单位元态射。它称为 A 的 Gelfand 变换。

若 \(x \in \mathrm{A}\)，则 Gelfand 变换 \(\mathcal{G}_{\mathrm{A}}(x)\)（即 \(x\) 的 Gelfand 变换）的像含于 \(\mathrm{Sp}_{\mathrm{A}}(x)\)。事实上，设 \(\chi \in \mathsf{X}(\mathrm{A})\)；由于 \(\chi(x - \chi(x)e) = 0\)，元素 \(x - \chi(x)e\) 不可逆。

设 B 是 K 上的一个交换幺代数，h 是从 A 到 B 的一个保单位元态射；于是 \(\mathsf{X}(h)\colon\mathsf{X}(\mathsf{B})\to\mathsf{X}(\mathsf{A})\) 定义了一个保单位元态射 \(h_{*}\colon\mathsf{K}^{\mathsf{X}(\mathsf{A})}\to\mathsf{K}^{\mathsf{X}(\mathsf{B})}\)，且图表：

\[
\begin{array}{l} \text {A} \xrightarrow {\mathcal {G} _ {\text {A}}} \text {K} ^ {\mathsf {X} (\text {A})} \\ \Big \downarrow h \qquad \qquad \qquad \Big \downarrow h _ {*} \\ \text {B} \xrightarrow {\mathcal {G} _ {\text {B}}} \text {K} ^ {\mathsf {X} (\text {B})} \end{array}\tag{3}
\]

是交换的。事实上，对每个 \(x \in A\) 与每个 \(\chi \in \mathsf{X}(\mathsf{B})\)，我们有：

\[
\begin{array}{r l} & {\mathcal {G} _ {\mathrm{B}} (h (x)) (\chi) = \chi (h (x)) = (\chi \circ h) (x)} \\ & {\qquad = (\mathsf {X} (h) (\chi)) (x)} \\ & {\qquad = \mathcal {G} _ {\mathrm{A}} (x) (\mathsf {X} (h) (\chi))} \\ & {\qquad = h _ {*} (\mathcal {G} _ {\mathrm{A}} (x)) (\chi).} \end{array}\tag{4}
\]

现在设 K 是一个拓扑域。这时我们给 \(\mathsf{X}(\mathsf{A})\) 赋予 A 上逐点收敛的拓扑（参见 EVT, III, p.~14, 例 1），而拓扑空间 \(\mathsf{X}(\mathsf{A})\) 称为 A 的特征标空间。于是 \(\mathsf{X}(\mathsf{A})\) 的拓扑是使诸函数 \(\mathcal{G}_{\mathrm{A}}(x)\)（\(x \in A\)）都连续的最粗拓扑，且映射 \(\chi \mapsto (\chi(a))_{a \in \mathrm{A}}\) 把空间 \(\mathsf{X}(\mathsf{A})\) 等同于 \(K^{A}\) 的一个子集。

当 K = R 或 C 时，这个拓扑是在 \(\mathsf{X}(\mathsf{A}) \subset \mathsf{A}^{*}\) 上由弱拓扑 \(\sigma(\mathsf{A}^{*}, \mathsf{A})\)（定义在 \(A^{*}\) 上）诱导的拓扑 (EVT, II, p.~45, 定义 2)；据此，我们也称它为 \(\mathsf{X}(\mathsf{A})\) 上的弱拓扑。

若 \(h\) 是从 A 到 B 的一个保单位元态射，则映射 \(\mathsf{X}(h)\colon \mathsf{X}(\mathsf{B})\to \mathsf{X}(\mathsf{A})\) 连续。若 \(h\) 是满射，则 \(\mathsf{X}(h)\) 的像，即在 \(h\) 的核上为零的 A 的特征标所成的集合，在 \(\mathsf{X}(\mathsf{A})\) 中是闭的；另一方面，\(\mathsf{X}(h)(\mathsf{X}(\mathsf{B}))\) 上由 \(\mathsf{X}(\mathsf{B})\) 的拓扑经由双射 \(\mathsf{X}(h)\) 诱导而得的拓扑，是 A 上逐点收敛的拓扑，也就是由 \(\mathsf{X}(\mathsf{A})\) 的拓扑诱导的拓扑。换言之，\(\mathsf{X}(h)\) 是从 \(\mathsf{X}(\mathsf{B})\) 到 \(\mathsf{X}(\mathsf{A})\) 的一个闭子集上的同胚。

若 \(A_{1},\ldots,A_{n}\) 是 K 上的交换幺代数，则空间 \(\mathsf{X}(\mathsf{A}_{1}\times\cdots\times\mathsf{A}_{n})\) 就这样等同于 \(\mathsf{X}(\mathsf{A}_{1}),\ldots,\mathsf{X}(\mathsf{A}_{n})\) 的拓扑和空间。类似地，\(\mathsf{X}(\mathsf{A}_{1}\otimes\cdots\otimes\mathsf{A}_{n})\) 等同于积拓扑空间 \(\mathsf{X}(\mathsf{A}_{1})\times\cdots\times\mathsf{X}(\mathsf{A}_{n})\)。

\subsection{无单位元代数的情形}

定义 6. --- 设 A 是 K 上的一个交换代数。A 的特征标是指从 A 到 K 的一个代数同态。

A 的特征标所成的集合将记作 X'(A)。

零映射是一个代数同态。若 A 有单位元 \(e\)，则从 A 到 K 中的每个非零代数态射都是保单位元的，也就是说，它是定义 4 意义下的 A 的一个幺特征标：事实上，对 \(\chi \in \mathsf{X}'(\mathsf{A})\)，\(\chi\) 非零必须且只需 \(\chi(e) = 1\)。

我们令 \(\mathsf{X}(\mathsf{A})=\mathsf{X}^{\prime}(\mathsf{A})-\{0\}\)；根据上文，这一记号与 A 为幺代数时引入的记号相容。

若 \(h\colon A\to B\) 是交换代数的一个态射，则映射 \(\chi\mapsto\chi\circ h\) 是一个映射 \(\mathsf{X}^{\prime}(h)\colon\mathsf{X}^{\prime}(\mathrm{B})\to\mathsf{X}^{\prime}(\mathrm{A})\)。它把 0 映到 0。若 \(k\colon B\to C\) 是交换代数的一个态射，则有 \(\mathsf{X}^{\prime}(k\circ h)=\mathsf{X}^{\prime}(h)\circ\mathsf{X}^{\prime}(k)\)。若 \(h\) 是满射，则 \(\mathsf{X}^{\prime}(h)\) 是从 \(\mathsf{X}^{\prime}(\mathrm{B})\) 到在 \(h\) 的核上为零的 A 的特征标所成集合上的一个双射。设 \(\mathrm{A}_{1},\ldots,\mathrm{A}_{n}\) 是交换代数，\(\mathrm{A}=\mathrm{A}_{1}\times\cdots\times\mathrm{A}_{n}\)，而 \(\pi_{i}\colon\mathrm{A}\to\mathrm{A}_{i}\) 是典范态射；则 \(\mathsf{X}^{\prime}(\pi_{i})\) 是从 \(\mathsf{X}^{\prime}(\mathrm{A}_{i})\) 到子集 \(\mathsf{X}_{i}^{\prime}\)（\(\mathsf{X}^{\prime}(\mathrm{A})\) 的一个子集）上的双射，这个子集即由 A 的在 \(\prod_{j\neq i}\mathrm{A}_{j}\) 上为零的特征标所成的集合；我们如第 6 小节那样看出，\(\mathsf{X}^{\prime}(\mathrm{A})\) 是

诸 \(X_{i}^{\prime}\) 的并；另一方面，有 \(X_{i}^{\prime} \cap X_{j}^{\prime} = \{0\}\)（对 \(i \neq j\)）；特别地，诸 \(X_{i}^{\prime} - \{0\}\) 构成 \(X^{\prime}(A) - \{0\} = X(A)\) 的一个划分。

对每个 \(x \in A\)，令 \(\mathcal{G}_{\mathrm{A}}^{\prime}(x)\)（或简作 \(\mathcal{G}^{\prime}(x)\)）为映射 \(\chi \mapsto \chi(x)\)（从 \(\mathsf{X}^{\prime}(\mathrm{A})\) 到 K）。映射 \(\mathcal{G}^{\prime}\) 是从 A 到代数 \(A_{1}\)（由在 0 处为零的映射 \(\mathsf{X}^{\prime}(\mathrm{A}) \to \mathrm{K}\) 所成的代数）中的一个态射。设 B 是一个交换代数，\(B_{1}\) 是由在 0 处为零的映射 \(\mathsf{X}^{\prime}(\mathrm{B}) \to \mathrm{K}\) 所成的代数，h 是从 A 到 B 的一个态射；则 \(\mathsf{X}^{\prime}(h)\) 定义了一个态射 \(h_{1}: A_{1} \to B_{1}\)，且我们有 \(h_{1} \circ \mathcal{G}_{A}^{\prime} = \mathcal{G}_{B}^{\prime} \circ h\)。用 \(\mathcal{G}_{\mathrm{A}}(x)\)（或简作 \(\mathcal{G}(x)\)）表示 \(\mathcal{G}_{\mathrm{A}}^{\prime}(x)\) 在 \(\mathrm{X}(\mathrm{A})\) 上的限制，并称它为 \(x\) 的 Gelfand 变换。

设 \(\widetilde{\mathsf{A}}\) 是由 A 添加单位元而得到的幺代数。通过限制，\(\widetilde{\mathsf{A}}\) 的每个特征标都定义 A 的一个特征标；反之，A 的每个特征标可唯一地延拓为 \(\widetilde{\mathsf{A}}\) 的特征标。这就定义了一个从 \(\mathsf{X}'(\mathsf{A})\) 到 \(\mathsf{X}(\widetilde{\mathsf{A}})\) 上的典范双射，据此把这两个集合视为同一。A 的特征标 0 与 \(\widetilde{\mathsf{A}}\) 的以 A 为核的唯一特征标视为同一。

若 \(x\in \mathrm{A}\) 且 \(\chi \in \mathsf{X}'(\mathrm{A})\)，则有 \(\chi (x)\in \mathrm{Sp}_{\widetilde{\mathrm{A}}}^{\sim}(x)\)，从而 \(\chi (x)\in \mathrm{Sp}_{\mathrm{A}}^{\prime}(x)\)。

引理 3. --- 映射 \(\chi \mapsto \mathrm{Ker}(\chi)\) 是从 \(\mathsf{X}(\mathsf{A})\) 到 A 的余维数为 1 的正则理想所成的集合上的一个双射。

回顾 (A, VIII, p.~426, 定义 1)，A 的理想 I 称为正则的，如果商代数 A/I 有单位元。

我们来证明该引理。一方面，\(\mathsf{X}(\mathsf{A})\) 与 \(\widetilde{A}\) 的在 A 上非零的特征标所成的集合视为同一。另一方面，由 A, VIII, p.~428, 命题 4，映射 \(I \mapsto A \cap I\) 是从 \(\widetilde{A}\) 的异于 A 的极大理想所成的集合到 A 的正则极大理想所成的集合上的一个双射。于是引理由第 6 小节的结果得出。

现在设 K 是一个拓扑域。这时我们给 \(\mathsf{X}^{\prime}(\mathsf{A})\) 赋予 A 上逐点收敛的拓扑；记号 \(\mathsf{X}^{\prime}(\mathsf{A})\) 此后表示这样得到的拓扑空间。当 K = R 或 C 时，我们也称它为弱拓扑。对每个 \(x \in A\)，函数 \(\mathcal{G}_{\mathrm{A}}^{\prime}(x)\)（定义在 \(\mathsf{X}^{\prime}(\mathsf{A})\) 上）连续。

若 h 是从 A 到 B 的一个态射，则映射 \(\mathsf{X}'(h)\colon\mathsf{X}'(\mathsf{B})\to\mathsf{X}'(\mathsf{A})\) 连续。若 h 是满射，则 \(\mathsf{X}'(h)\) 是从 \(\mathsf{X}'(\mathsf{B})\) 到其像上的一个同胚，且这个像在 \(\mathsf{X}'(\mathsf{A})\) 中是闭的。

设 \(A = A_{1} \times \cdots \times A_{n}\)；沿用上面的记号，\(\mathsf{X}'(\pi_{i})\) 是从 \(\mathsf{X}'(\mathrm{A}_{i})\) 到 \(X_{i}'\) 上的一个同胚，且 \(X_{i}'\) 在 \(\mathsf{X}'(\mathrm{A})\) 中是闭的。因此 \(X_{i}' - \{0\}\) 在 \(\mathsf{X}'(\mathrm{A})\) 中是开的；诸 \(\mathsf{X}'(\pi_{i})\) 定义了从诸 \(\mathsf{X}'(\mathrm{A}_{i})\) 的和空间 S 到 \(\mathsf{X}'(\mathrm{A})\) 上的一个连续映射，并且可以立即验证，点 \(0 \in \mathsf{X}'(\mathrm{A}_{1}), \ldots, 0 \in \mathsf{X}'(\mathrm{A}_{n})\) 的邻域之并的像是 \(0 \in \mathsf{X}'(\mathrm{A})\) 的一个邻域；由此得出，\(\mathsf{X}'(\mathrm{A})\) 典范地等同于 S 的一个商空间。特别地，空间 \(\mathsf{X}(\mathrm{A})\) 与诸 \(\mathsf{X}(\mathrm{A}_{i})\) 的和空间视为同一。

从 \(\mathsf{X}^{\prime}(\mathsf{A})\) 到 \(\mathsf{X}(\tilde{\mathsf{A}})\) 的典范双射是一个同胚。设 B 是 K 上的一个幺代数，\(B^{\prime}\) 是其基础代数；则空间 \(\mathsf{X}(\mathsf{B})\) 与子空间 \(\mathsf{X}(\mathsf{B}^{\prime})\)（\(\mathsf{X}^{\prime}(\mathsf{B}^{\prime})\) 的子空间）视为同一。

\subsection{本原理想}

设 A 是 K 上的一个代数，E 是 K 上的一个向量空间。A 在 E 中的一个表示是指从 A 到 E 的自同态所成的代数 \(\mathcal{L}(\mathrm{E})\) 的一个态射。单射的表示称为忠实的。设 \(\pi_1\) 与 \(\pi_2\) 是 A 在空间 \(\mathrm{E}_1\)、\(\mathrm{E}_2\) 中的表示。从 \(\pi_1\) 到 \(\pi_2\) 的一个态射是指一个 K-线性映射 \(u: \mathrm{E}_1 \to \mathrm{E}_2\)，它使得 \(u(\pi_1(a)x) = \pi_2(a)u(x)\) 对一切 \(a \in \mathrm{A}\) 与 \(x \in \mathrm{E}_1\) 成立。若存在一个从 \(\pi_1\) 到 \(\pi_2\) 的态射，它是向量空间的一个同构，则称这些表示是等价的。其逆于是是从 \(\pi_2\) 到 \(\pi_1\) 的一个态射。A 在 E 中的表示 \(\pi\) 称为不可约的，如果 \(\mathrm{E} \neq \{0\}\)，并且 E 的在 \(\pi(\mathrm{A})\) 下稳定的向量子空间只有 \(\{0\}\) 与 E。

例. --- 从 A 到 \(\mathcal{L}(\mathrm{E})\) 的零映射是一个表示，称为 A 的平凡表示。它不可约必须且只需 E 的维数为 1。

引理 4. --- 设 π 是 A 在 E 中的一个不可约的非平凡表示。对 E 的每个非零元素 ξ，有 π(A)ξ = E。

子空间 \(\pi(A)\xi\)（\(E\) 的子空间）在 \(\pi(A)\) 下稳定。假定它为零。那么非零子空间 \(K\xi\)（\(E\) 的非零子空间）就会在 \(\pi(A)\) 下稳定，从而等于 \(E\)；但这将蕴涵 \(\pi\) 是零表示。因此我们有 \(\pi(A)\xi = E\)。

设 \(\pi\) 是 A 在 E 中的一个不可约的非平凡表示。根据该引理，\(\xi\) 在 A 中的零化子 R 是 A 的一个正则左理想 (A, VIII, p.~425, 第 1 小节)，且表示 \(\pi\) 等价于由 A-伪模 A/R 所定义的表示。由于 \(\pi\) 不可约，理想 R 是一个极大正则左理想。

定义 7. --- 设 A 是 K 上的一个代数。A 的本原理想是指 A 的一个非平凡不可约表示的核。

若 A 交换，则 A 的本原理想就是 A 的正则极大理想。事实上，A 的非平凡不可约表示，除了相差一个等价外，就是由 A-伪模 A/R 所定义的表示 \(\pi_{R}\)，其中 R 是 A 的一个正则极大理想。\(\pi_{R}\) 的核包含 R。它甚至等于 R，因为由 A, VIII, p.~426, 命题 2，A 的交换性蕴涵 A/R 是一个域。因此 \(\text{Ker}(\pi_{\text{R}})\) 是正则极大的。

引理 5. --- 设 π 是 A 在 K 上一个向量空间 E 中的一个不可约表示。

\begin{enumerate}
\def\labelenumi{\alph{enumi})}
\item
  设 I 是 A 的一个双边理想。若 \(\pi(\mathrm{I}) \neq \{0\}\)，则 \(\pi|\mathrm{I}\) 不可约；
\item
  设 \(I_{1}\) 与 \(I_{2}\) 是 A 的双边理想，使得 \(\pi(\mathrm{I}_{1}) \neq 0\) 且 \(\pi(\mathrm{I}_{2}) \neq 0\)。则 \(\pi(\mathrm{I}_{1}\mathrm{I}_{2}) \neq 0\)。
\end{enumerate}

E 中被 \(\pi(\mathrm{I})\) 零化的元素所成的集合在 \(\pi(\mathrm{A})\) 下稳定且异于 E，故等于 0。于是，若 \(\xi\) 是 E 的一个非零元素，则有 \(\pi(\mathrm{I})\xi \neq 0\)；由于 \(\pi(\mathrm{I})\xi\) 在 \(\pi(\mathrm{A})\) 下稳定，有 \(\pi(\mathrm{I})\xi = \mathrm{E}\)，这就证明了 \(a\) )。另一方面，上文证明了 \(\pi(\mathrm{I}_2)\mathrm{E} = \mathrm{E}\)，\(\pi(\mathrm{I}_1)\pi(\mathrm{I}_2)\mathrm{E} = \mathrm{E}\)，从而 \(\pi(\mathrm{I}_1\mathrm{I}_2) \neq 0\)，由此得 \(b\) )。

引理 6. --- 设 I \(_{1}\) 与 I \(_{2}\) 是 A 的双边理想，I 是 A 的一个本原理想。若 I 包含 I \(_{1}\) I \(_{2}\)（特别地，若 I 包含 I \(_{1}\) ∩ I \(_{2}\)），则 I 包含 I \(_{1}\) 或 I \(_{2}\)。

设 \(\pi\) 是以 I 为核的一个不可约表示。若 I \(\not\supseteq\) I \(_1\) 且 I \(\not\supseteq\) I \(_2\)，则引理 5 的 \(b\) ) 证明 \(\pi(\text{I}_1\text{I}_2) \neq 0\)，从而 I \(\not\supseteq\) I \(_1\) I \(_2\)。

引理 7. --- 假定 A 有单位元。设 I 是 A 的一个极大双边理想。则 I 是一个本原理想。

存在 A 的一个包含 I 的极大左理想 R (A, I, p.~99, 定理 1)。设 \(\pi\) 是 A 在 A/R 中的典范表示，它不可约且非零。由于 IA \(\subset\) R，\(\pi\) 的核 I′ 包含 I，故 I′ = I，从而 I 是本原的。

设 J(A) 是 A 的本原理想所成的集合。对 A 的任何子集 M，我们用 V(M) 表示包含 M 的 A 的本原理想所成的集合；若 I 是由 M 生成的 A 的双边理想，则 V(M) = V(I)。若 M 化为单个元素 x，我们就写 V(x) 以代替 V(\{x\})。

映射 M \(\mapsto\) V(M) 关于包含关系是递减的。我们有：

\[
\begin{array}{r l} (5) & \mathrm{V} (\varnothing) = \mathrm{J} (\mathrm{A}), \qquad \mathrm{V} (\mathrm{A}) = \varnothing \end{array}
\]

\[
\begin{array}{r l} (6) & \mathrm{V} \left(\bigcup_ {i \in \mathrm{I}} \mathrm{M} _ {i}\right) = \mathrm{V} \left(\sum_ {i \in \mathrm{I}} \mathrm{M} _ {i}\right) = \bigcap_ {i \in \mathrm{I}} \mathrm{V} (\mathrm{M} _ {i}) \end{array}
\]

对 A 的子集的任何族 \((\mathrm{M}_{i})_{i\in\mathrm{I}}\) 成立。另一方面，由引理 6，

\[
\begin{array}{r l} (7) & \mathrm{V} (\mathrm{I} _ {1} \cap \mathrm{I} _ {2}) = \mathrm{V} (\mathrm{I} _ {1} \mathrm{I} _ {2}) = \mathrm{V} (\mathrm{I} _ {1}) \cup \mathrm{V} (\mathrm{I} _ {2}) \end{array}
\]

对 A 的任意双边理想 \(I_{1}\)、\(I_{2}\) 成立。公式 (5) 至 (7) 表明，J(A) 的子集 V(M) 是一个称为 J(A) 上 Jacobson 拓扑的拓扑的闭子集。

设 T 是 J(A) 的一个子集，并设 \(\Upsilon(T)\) 是 T 的诸元素之交，于是 \(\Upsilon(T)\) 是 A 的一个双边理想。则 T 在 J(A) 中的闭包是包含 T 的最小闭子集，即 V( \(\Upsilon(T)\) )。特别地，T 是闭的必须且只需 \(T = V(\Upsilon(T))\)。

命题 2. --- 设 I \(_{1}\) 与 I \(_{2}\) 是 J(A) 的两个不同的点。则这两点中的一个不附着于另一个。

事实上，例如我们有 \(I_{1} \not\subset I_{2}\)。集合 \(V(I_{1})\)，它由 \(I \in J(A)\)（使得 \(I_{1} \subset I\)）组成，在 \(J(A)\) 中是闭的，它包含 \(I_{1}\) 但不包含 \(I_{2}\)。

命题 3. --- 设 I ∈ J(A)。要使 \{I\} 在 J(A) 中是闭的，必须且只需 I 是一个本原极大理想。

事实上，\{I\} 的闭包由包含 I 的 A 的本原理想组成。

关系

« \(\pi_{1},\pi_{2}\) 是 A 的表示，且它们同构。

在 \(\pi_1\) 与 \(\pi_2\) 之间是一个等价关系。对 A 的每个表示 \(\pi\)，我们用 \(\operatorname{cl}(\pi)\) 表示 \(\pi\) 的等价类，它因此是一个与 \(\pi\) 同构的 A 的表示，使得两个表示 \(\pi_1\) 与 \(\pi_2\) 同构必须且只需 \(\operatorname{cl}(\pi_1) = \operatorname{cl}(\pi_2)\)。我们说 \(\operatorname{cl}(\pi)\) 是 \(\pi\) 的类。

设 c 是 A 的基数。设 \(\pi\) 是 A 在一个 K-向量空间 E 中的一个非零不可约表示。设 \(\xi\) 是 E 的一个非零元素。由于 \(\pi(A)\xi = E\) (引理 4)，E 的维数 \(\leqslant\mathfrak{c}\) (A, II, p.~97, 推论)。关系

“λ 是 A 的不可约表示的一个类

在一个维数 \(\leqslant\mathfrak{c}\gg\)

在 \(\lambda\) 中是可汇集的 (E, II, p.~3)。事实上，每个维数 \(\leqslant \mathfrak{c}\) 的向量空间都同构于一个空间 \(\mathrm{K}^{\mathrm{B}}\)，其中 B 是 A 的一个子集 (A, II, p.~25, 定义 10)，而断言则由 E, II, p.~47 得出。

我们用 \(\widehat{\mathrm{A}}\) 表示 A 的非平凡不可约表示的类所成的集合。根据上文，对 A 的每个非平凡不可约表示 \(\pi\)，存在唯一的一个表示 \(\widehat{\pi} \in \widehat{\mathrm{A}}\)，它与 \(\pi\) 同构。

从 \(\hat{\mathrm{A}}\) 到 \(\mathrm{J(A)}\) 中的、把 \(\pi\) 对应到其核的映射是满射。若 A 交换，则由本原理想是正则极大理想这一事实推出，这个映射是一个双射。

我们给 \(\hat{A}\) 赋予在映射 \(\hat{A} \rightarrow J(A)\) 下的原像拓扑。

若 A 有单位元，则空间 J(A) 与 \(\hat{\mathrm{A}}\) 是拟紧的。

只须对 J(A) 给出证明。设 \((\mathrm{T}_{j})\) 是 J(A) 的闭子集的一个族，其交为空。若和 \(\sum_{j}\Upsilon(\mathrm{T}_{j})\) 异于 A，则这个和将包含于一个极大双边理想 I。理想 I 将是本原的 (引理 7)；由于每个 \(T_{j}\) 是闭的，故等于 \(\mathrm{V}(\Upsilon(\mathrm{T}_{j}))\)，我们就会有 \(I\in T_{j}\) 对一切 j 成立，这与假设矛盾。于是我们有 \(\sum_{j}\Upsilon(\mathrm{T}_{j})=\mathrm{A}\)，因而可以写 \(1=x_{1}+\cdots+x_{n}\)，其中 \(n\geqslant1\)，且对一切 i 有 \(x_{i}\in\Upsilon(\mathrm{T}_{j_{i}})\)。这蕴涵 \(\Upsilon(\mathrm{T}_{j_{1}})+\cdots+\Upsilon(\mathrm{T}_{j_{n}})=\mathrm{A}\)，从而 \(T_{j_{1}}\cap\cdots\cap T_{j_{n}}=\varnothing\)。

假定代数 A 是交换的且有单位元。J(A) 上的 Jacobson 拓扑是由 A 的素谱的 Zariski 拓扑在 J(A) 上诱导的拓扑 (AC, II, 定义 4, p.~125)。

假定 A 交换且 K 是一个拓扑域。从 K 到 \(\mathcal{L}(\mathrm{K})\) 上的典范同构使我们可以把 \(\mathsf{X}(\mathsf{A})\) 的一个元素等同于 A 在向量空间 K 中的一个表示，这就定义了一个从 \(\mathsf{X}(\mathsf{A})\) 到 \(\widehat{\mathsf{A}}\) 中的单射。因此可以把 \(\mathsf{X}(\mathsf{A})\) 与 \(\widehat{\mathsf{A}}\) 的一个子集视为同一。

命题 5. — 在 \(\mathsf{X}(\mathsf{A})\) 上由 \(\widehat{\mathsf{A}}\) 的拓扑诱导的拓扑粗于 \(\mathsf{X}(\mathsf{A})\) 的拓扑。

事实上，设 T 是 \(\widehat{\mathbf{A}}\) 的一个闭子集。则 T 是由其核包含 A 的一个子集 M 的那些 \(\pi \in \widehat{\mathbf{A}}\) 所成的集合。因此 \(\mathrm{T} \cap \mathrm{X}(\mathrm{A})\) 是由在 M 上为零的那些 \(\chi \in \mathrm{X}(\mathrm{A})\) 所成的集合，即 \(\mathrm{X}(\mathrm{A})\) 的一个闭子集。由此得命题。

一般说来，\(\mathsf{X}(\mathsf{A})\) 的拓扑与由 \(\widehat{\mathsf{A}}\) 的拓扑诱导的拓扑并不重合 (参见 I, p.~193, 习题 6, c))。

\section{赋范代数}

在本节中，K 表示域 R 或 C 之一。

\subsection{概论}

回顾 (参见 TG, IX, p.~37, 定义 9)，赋范代数是指 K 上的一个代数 A，它带有一个范数 \(x \mapsto \|x\|\)，使得：

\[
\| x y \| \leqslant \| x \| \| y \|\tag{1}
\]

对一切 \(x, y \in A\) 成立。若 A 完备，则称 A 是一个 Banach 代数。

设 A 是 K 上的一个完备可赋范代数。A 的拓扑可由一个满足 (1) 的范数定义 (参见 TG, IX, p.~38)。当 A 带有由这样一个范数所定义的赋范代数结构时，我们也说代数 A 是一个 Banach 代数。

我们还回顾下列事实 (参见 TG, IX, p.~38--39)：

\begin{enumerate}
\def\labelenumi{(\arabic{enumi})}
\item
  若 A 是一个非零赋范代数且有单位元 e，则 \(\|e\|\geqslant1\)。
\item
  设 A 是一个赋范 K-代数。A 的反代数，带上同一个范数，是一个赋范代数。A 的完备化，带上由 A 的范数连续延拓而得的范数，是一个 K-Banach 代数。A 的每个子代数，带上诱导范数，是一个赋范代数。若 I 是 A 的一个闭双边理想，则代数 A/I，带上由 \(\|\dot{x}\| = \inf_{x \in \dot{x}} \|x\|\)（对一切 \(\dot{x} \in A/I\)）所定义的范数，是一个赋范代数。
\item
  设 \((\mathrm{A}_{i})_{i\in\mathrm{I}}\) 是一族赋范代数，B 是诸代数 \(A_{i}\) 的积代数。由 B 的元素 \((x_{i})_{i\in\mathrm{I}}\)（满足 \(\|(\boldsymbol{x}_{i})\|=\sup_{i\in\mathrm{I}}\|\boldsymbol{x}_{i}\|<+\infty\)）所成的子代数是一个赋范代数，称为诸代数 \((\mathrm{A}_{i})_{i\in\mathrm{I}}\) 的积赋范代数。若对每个 i，\(A_{i}\) 都是 Banach 代数，则 B 是 Banach 代数。
\end{enumerate}

设 A 是一个赋范代数。设 \((y_{i})_{i\in\mathrm{I}}\) 是 A 的元素的一个族；包含诸元素 \(y_{i}\) 的最小闭子代数 B 称为由诸 \(y_{i}\) 生成的 A 的闭子代数；若 B = A，就说诸 \(y_{i}\) 拓扑地生成赋范代数 A，或者说族 \((y_{i})_{i\in\mathrm{I}}\) 是赋范代数 A 的一个拓扑生成系。

类似地，若 A 是一个赋范幺代数，则包含诸元素 \(y_{i}\) 的最小闭幺子代数称为由诸 \(y_{i}\) 生成的 A 的闭幺子代数。若它等于 A，就说诸 \(y_{i}\) 拓扑地生成赋范幺代数 A。

设 A 是一个赋范 K-代数。用 \(\tilde{\mathbf{A}}\) 表示由 A 添加单位元而得到的代数。在 \(\tilde{\mathbf{A}}\) 上，令 \(\| (\lambda ,x)\| = |\lambda | + \| x\|\)（对一切 \(\lambda \in \mathbf{K}\) 与每个 \(x\in \mathbf{A}\)），这样就定义了一个范数。我们有

\[
\begin{array}{r l} & {\| (\lambda , x) (\mu , y) \| = | \lambda \mu | + \| x y + \mu x + \lambda y \|} \\ & {\qquad \leqslant | \lambda | | \mu | + \| x \| \| y \| + | \mu | \| x \| + | \lambda | \| y \|} \\ & {\qquad = \| (\lambda , x) \| \| (\mu , y) \|.} \end{array}
\]

于是 \(\widetilde{\mathsf{A}}\) 成为一个赋范代数，称为由 A 添加单位元而得到的赋范幺代数。代数 A 与闭双边理想 \(\{0\} \times \mathbf{A}\)（\(\widetilde{\mathsf{A}}\) 的闭双边理想）视为同一。

定义 1. --- 设 A 是一个赋范代数。对 \(a \in \mathrm{A}\)，用 \(\gamma_{a}\) 与 \(\delta_{a}\) 表示从 A 到 A 的映射 \(x \mapsto ax\) 与 \(x \mapsto xa\)。映射 \(\gamma: a \mapsto \gamma_{a}\) 是 A 在 A 中的一个表示，称为 A 的左正则表示。映射 \(\delta: a \mapsto \delta_{a}\) 是 A 的反代数在 A 中的一个表示，称为 A 的右正则表示。

引理 1. --- 设 A 是一个赋范代数。A 的左正则表示与 A 的右正则表示都是连续的，且范数 \(\leqslant 1\)。

若代数 A 是幺的，则映射 \(x \mapsto \|\gamma_{x}\|\) 与 \(x \mapsto \|\delta_{x}\|\) 是 A 上的范数，它们定义 A 的拓扑。它们满足不等式 (1)。

若赋范幺代数 A 非零，也就是说，若 \(e \neq 0\)，则我们还有 \(\|\gamma_{e}\| = \|\delta_{e}\| = 1\)。

立即有 \(\|\gamma_{x}\|\leqslant\|x\|\) 与 \(\|\delta_{x}\|\leqslant\|x\|\)。

若 A 有单位元 \(e\)，则我们有 \(x = \gamma_x e = \delta_x e\)，因而：

\[
\| x \| \leqslant \| \boldsymbol {\gamma} _ {x} \| \cdot \| e \| \quad \| x \| \leqslant \| \boldsymbol {\delta} _ {x} \| \cdot \| e \|.\tag{2}
\]

这时，\(x \mapsto \|\gamma_{x}\|\) 与 \(x \mapsto \|\delta_{x}\|\) 是与 A 的范数等价的范数。它们满足 (1)，因为 \(\gamma\) 与 \(\delta\) 是表示。

最后一个断言由 \(\gamma_{e} = \delta_{e} = Id_{A}\) 这一事实得出。

\subsection{例}

\begin{enumerate}
\def\labelenumi{\arabic{enumi})}
\item
  设 E 是一个赋范空间。给 E 带上由 ab = 0（对 E 中所有的 a 与 b）定义的乘法；这样就在 E 上定义了一个赋范代数结构。
\item
  设 X 是一个集合。我们用 \(\mathcal{B}(X;K)\) 表示 X 上取值于 K 的有界函数所成的赋范代数，带上范数

\[
\| f \| = \sup _ {x \in X} | f (x) |
\]

(TG, X, p.~22)。它是 K 上的一个幺 Banach 代数 (TG, X, p.~21, 推论 1)。它是交换的。设 \(f\) 是 \(\mathcal{B}(\mathrm{X};\mathrm{K})\) 的一个元素。则 \(f\) 在 \(\mathcal{B}(\mathrm{X};\mathrm{K})\) 中可逆必须且只需

\[
\inf _ {x \in X} | f (x) | > 0.
\]

于是 \(f\) 的谱是由使下式成立的 \(\lambda\in K\) 所成的集合：

\[
\inf _ {x \in \mathrm{X}} | f (x) - \lambda | = 0,
\]

也就是值集 \(f(\mathrm{X})\)（\(f\) 的值集）在 K 中的闭包。

\item
  设 X 是一个拓扑空间。我们用 \(\mathcal{C}_{b}(\mathrm{X};\mathrm{K})\) 表示 \(\mathcal{B}(\mathrm{X};\mathrm{K})\) 的由 X 上取值于 K 的连续且有界的函数所成的幺子代数，用 \(\mathcal{C}_{0}(\mathrm{X};\mathrm{K})\) 表示 \(\mathcal{C}_{b}(\mathrm{X};\mathrm{K})\) 的由在无穷远处趋于 0 的函数所成的子代数 (参见 INT, III, §1, 第 2 小节及 TG, X, p.~40, 推论 2)。我们回顾，\(\mathcal{H}(\mathrm{X};\mathrm{K})\) 表示 \(\mathcal{C}_{b}(\mathrm{X};\mathrm{K})\) 的由具有紧支集的连续函数所成的子代数。
代数 \(\mathcal{C}_{b}(\mathrm{X};\mathrm{K})\) 与 \(\mathcal{C}_{0}(\mathrm{X};\mathrm{K})\) 是 K 上的交换 Banach 代数；事实上，它们是 \(\mathcal{B}(\mathrm{X};\mathrm{K})\) 的闭子空间 (TG, X, p.~21, 推论 2 及 INT, III, §1, 第 2 小节)。

由于一个处处非零的连续函数的逆是连续的，子代数 \(\mathcal{C}_{b}(\mathrm{X};\mathrm{K})\) 是 \(\mathcal{B}(\mathrm{X};\mathrm{K})\) 的一个全子代数。特别地，\(\mathcal{C}_{b}(\mathrm{X};\mathrm{K})\) 的元素 f 的谱等于 \(\overline{f(\mathrm{X})}\)。

若 X 是离散的，则有 \(\mathcal{C}_b(\mathrm{X};\mathrm{K}) = \mathcal{B}(\mathrm{X};\mathrm{K})\)。

若 X 是紧的，则有 \(\mathcal{C}_{b}(\mathrm{X};\mathrm{K}) = \mathcal{H}(\mathrm{X};\mathrm{K}) = \mathcal{C}(\mathrm{X};\mathrm{K})\)，即 X 上取值于 K 的连续函数所成的赋范幺代数 \(\mathcal{C}(\mathrm{X};\mathrm{K})\)。这时，元素 \(f \in \mathcal{C}(\mathrm{X};\mathrm{K})\) 可逆必须且只需 f 不取值 0，且 f 的谱等于 f 的值集 \(f(\mathrm{X})\)。

此后假定 X 不是紧的。于是代数 \(\mathcal{C}_{0}(X;K)\) 不是幺的；由它添加单位元而得到的代数与子代数 \(K \cdot 1 \oplus \mathcal{C}_{0}(X;K)\)（\(\mathcal{C}(X;K)\) 的由在无穷远处有极限的函数所成的子代数）视为同一。这个子代数的一个元素可逆必须且只需它不为零且它在无穷远处的极限非零。由此得出，对每个 \(f \in \mathcal{C}_{0}(X;K)\)，f 的谱等于 \(f(X) \cup \{0\}\)。

\(f \in \mathcal{K}(\mathrm{X};\mathrm{K})\) 的谱等于 \(f(\mathrm{X})\)（事实上，若 \(\mathrm{X}\) 不是紧的，则 0 属于 \(f(\mathrm{X})\)）。

\item
  设 \(n \geqslant 0\) 是一个整数。设 \(A_{n}\) 是由函数 \(f: [0,1] \to K\)（在 \([0,1]\) 上有直到 \(n\) 阶的连续导数）所成的代数，带上范数
\[
\| f \| = \sum_ {k = 0} ^ {n} \frac {1}{k !} \sup _ {0 \leqslant t \leqslant 1} | f ^ {(k)} (t) |.
\]

若 \(f,g\in \mathrm{A}_n\)，则有

\[
\begin{array}{l} \| f g \| = \sum_ {k = 0} ^ {n} \frac {1}{k !} \sup | (f g) ^ {(k)} (t) | = \sum_ {k = 0} ^ {n} \frac {1}{k !} \sup \Bigl | \sum_ {s = 0} ^ {k} \binom {k} {s} f ^ {(s)} (t) g ^ {(k - s)} (t) \Bigr | \\ \leqslant \sum_ {k = 0} ^ {n} \sum_ {s = 0} ^ {k} \frac {1}{s ! (k - s) !} \sup _ {0 \leqslant t \leqslant 1} | f ^ {(s)} (t) | \sup _ {0 \leqslant t \leqslant 1} | g ^ {(k - s)} (t) | = \| f \|   \| g \|, \end{array}
\]

这是由 Leibniz 公式 (FVR, I, p.~28, 命题 2) 得出的，因而 \(A_{n}\) 是一个交换的幺 Banach 代数。

\item
  设 E 是一个 Banach 空间，其范数记作 p。E 的连续自同态所成的代数 \(\mathcal{L}(\mathrm{E})\)，带上范数
\[
\| u \| = \sup _ {p (x) \leqslant 1} p (u (x))
\]

是一个幺 Banach 代数 (EVT, III, p.~14 及 p.~24, 推论 2; TG, X, p.~23, 公式 (3))。

\item
  设 G 是一个局部紧群。用 e 表示它的单位元。设 \(\mathcal{M}^{1}(G)\) 是 G 上的有界复测度所成的 Banach 空间 (INT, III, p.~57)。卷积积 (INT, VIII, p.~120, 定义 1) 给 \(\mathcal{M}^{1}(G)\) 带上一个复 Banach 代数结构 (INT, VIII, §3, 第 1 小节, 命题 2)，其单位元是由置于点 e 的单位质量所定义的测度 \(\varepsilon_{e}\)。若 G 交换，这个 Banach 代数是交换的。具有紧支集的测度所成的空间 \(\mathcal{C}'(G)\) 是 \(\mathcal{M}^{1}(G)\) 的一个子代数 (同前)。
\item
  设 G 是一个带有 Haar 测度 \(\mu\) 的局部紧群。则 \(\mathrm{L}_{\mathrm{K}}^{1}(\mathrm{G},\mu)\) 关于卷积积是一个 Banach 代数 (INT, VIII, 命题 12, p.~166)。若 K = C，则由 \(f \mapsto f \mu\) 定义的映射使我们可以把 \(\mathrm{L}_{\mathbf{C}}^{1}(\mathrm{G},\mu)\) 与 Banach 代数 \(\mathcal{M}^{1}(\mathrm{G})\) 的一个子代数视为同一。
若 G 交换，则 Banach 代数 \(\mathrm{L}_{\mathrm{K}}^{1}(\mathrm{G},\mu)\) 是交换的。

\item
  在例 7 中取 G = Z 及 K = C。则 \(\mathrm{L}_{\mathbf{C}}^{1}(\mathbf{Z})\) 是由序列 \((x_{n})_{n\in\mathbf{Z}}\)（满足 \(\sum_{n}|x_{n}|<+\infty\)）所成的复交换 Banach 代数，元素 \((x_{n})\) 与 \((y_{n})\) 的乘积是 \((z_{n})\)，其中
\[
z _ {n} = \sum_ {k \in {\bf Z}} x _ {k} y _ {n - k}
\]

而范数为 \(\|(x_{n})\|=\sum_{n}|x_{n}|\)。这个代数以序列 \(\varepsilon=(\varepsilon_{n})\) 为单位元，其中 \(\varepsilon_{0}=1\)，且 \(\varepsilon_{n}=0\)（对 \(n\neq0\)）。

用 U 表示 C 中的单位圆。若 \(x = (c_{n})\) 是 A 的一个元素，用 \(\varphi(x)\) 表示 U 上的连续函数，它在 \(e^{it}\) 处的值为

\[
\varphi (x) (e ^ {i t}) = \sum_ {n \in {\bf Z}} c _ {n} e ^ {i n t}.
\]

可以验证，\(\varphi\) 是从 \(\mathrm{L}_{\mathbf{C}}^{1}(\mathbf{Z})\) 到 \(\mathbf{U}\) 上连续函数的一个代数 A 上的态射，A 中的乘法就是通常的乘法。把等式逐项积分

\[
\left(\sum_ {m \in {\bf Z}} c _ {m} e ^ {i m t}\right) \cdot e ^ {- i n t} = \varphi (x) (e ^ {i t}) \cdot e ^ {- i n t},
\]

就得到

\[
c _ {n} = \frac {1}{2 \pi} \int_ {0} ^ {1} \varphi (x) (e ^ {i t}) e ^ {- i n t} d t.
\]

特别地，由此得出态射 \(\varphi\) 是单射。代数 A 带上由 \(\mathrm{L}_{\mathbf{C}}^{1}(\mathbf{Z})\) 的范数经由 \(\varphi\) 诱导的范数，称为绝对收敛 Fourier 级数的 Banach 代数。它以函数 \(1 = \varphi(\varepsilon)\) 为单位元。

\item
  设 \(\Delta\) 是由复数 \(z\)（满足 \(|z| \leqslant 1\)）所成的圆盘。\(\Delta\) 上连续且在 \(\Delta\) 的内部全纯的函数所成的代数 (VAR, R1, p.~26, 3.2.1) 带上范数 \(\|f\| = \sup_{z \in \Delta} |f(z)|\)。则 A 是一个交换的幺 Banach 代数。
\end{enumerate}

\subsection{谱半径}

引理 2 (Fekete 引理). --- 设 \((a_{n})_{n\geqslant1}\) 是一个实数序列。假定

\[
a _ {n + m} \leqslant a _ {n} + a _ {m}
\]

对一切 \(n \geqslant 1\) 与每个 \(m \geqslant 1\) 成立。则序列 \((a_n / n)_{n \geqslant 1}\) 收敛且满足

\[
\lim _ {n \to + \infty} \frac {a _ {n}}{n} = \inf _ {n \geqslant 1} \frac {a _ {n}}{n}.
\]

令 \(a_0 = 0\)；不等式 \(a_{n+m} \leqslant a_n + a_m\) 对每个 \(n \geqslant 0\) 与每个 \(m \geqslant 0\) 仍然成立。固定一个整数 \(m \geqslant 1\)。对每个整数 \(n \geqslant 1\)，用 \(q(n)\) 与 \(r(n)\) 表示使 \(n = q(n)m + r(n)\) 且 \(0 \leqslant r(n) < m\) 的整数 (E, III, p.~39, 定理 1)。于是假设蕴涵

\[
\frac {a _ {n}}{n} \leqslant \frac {a _ {q (n) m}}{n} + \frac {a _ {r (n)}}{n} \leqslant \frac {q (n) a _ {m}}{n} + \frac {a _ {r (n)}}{n} \leqslant \frac {q (n)}{n} a _ {m} + \frac {m}{n}.
\]

令 \(n\) 趋于 \(+\infty\)，就推出 \(\limsup_{n}(a_{n}/n) \leqslant a_{m}/m\)，因为 \(q(n)/n \to 1/m\)。由于这对每个 \(m \geqslant 1\) 都成立，因此我们有

\[
\limsup _ {n \to + \infty} \frac {a _ {n}}{n} \leqslant \inf _ {m \geqslant 1} \frac {a _ {m}}{m} \leqslant \liminf _ {n \to + \infty} \frac {a _ {n}}{n}.
\]

这些不等式证明了序列 \((a_{n}/n)_{n\geqslant1}\) 的收敛性以及公式 \(\lim a_{n}/n = \inf_{n\geqslant1} a_{n}/n\)。

命题 1. --- 设 A 是一个赋范代数。对每个 \(x \in \mathrm{A}\)，序列 \((\|x^n\|^{1/n})_{n \geqslant 1}\) 收敛，且其极限 \(\varrho(x)\) 等于 \(\inf_{n \geqslant 1} \|x^n\|^{1/n}\)。而且，对定义 A 的拓扑的任何范数 \(x \mapsto \|x\|_1\)，也有

\[
\varrho (x) = \lim _ {n \to + \infty} \| x ^ {n} \| _ {1} ^ {1 / n} = \inf _ {n \geqslant 1} \| x ^ {n} \| _ {1} ^ {1 / n}.
\]

若 \(x\) 是幂零的，则有 \(\| x^n\| _1^{1 / n} = 0\)（对每个充分大的整数 \(n\) 以及定义 A 的拓扑的每个范数 \(x\mapsto \| x\| _1\)）。

现在设 \(x\) 不是幂零的，并令 \(\alpha_{n} = \| x^{n}\|\)。有 \(\alpha_{n} > 0\)（对每个整数 \(n \geqslant 1\)），且由 (1)，有 \(\alpha_{n + m} \leqslant \alpha_{n}\alpha_{m}\)（对一切 \(n, m \in \mathbf{N}\)）。把引理 2 应用于序列 \(a_n = \log(\alpha_n)\)，就证明了极限 \(\varrho(x)\) 的存在性以及公式 \(\varrho(x) = \inf_{n > 0} \alpha_n^{1/n}\)。

设 \(x \mapsto \|x\|_1\) 是定义 A 的拓扑的一个范数。存在实数 \(a > 0\) 与 \(b > 0\)，使得

\[
a \| x \| \leqslant \| x \| _ {1} \leqslant b \| x \|
\]

对一切 \(x \in A\) 成立 (EVT, II, p.~7, 推论 2)。由此得

\[
a ^ {1 / n} \| x ^ {n} \| ^ {1 / n} \leqslant \| x ^ {n} \| _ {1} ^ {1 / n} \leqslant b ^ {1 / n} \| x ^ {n} \| ^ {1 / n},
\]

对一切 \(n \geqslant 1\) 成立；因而，取极限或取下确界，就得到等式

\[
\varrho (x) = \lim _ {n \to + \infty} \| x ^ {n} \| _ {1} ^ {1 / n} = \inf _ {n > 0} \| x ^ {n} \| _ {1} ^ {1 / n}.
\]

定义 2. --- 对赋范代数 A 的每个元素 x，实数

\[
\varrho (x) = \lim _ {n \to \infty} \| x ^ {n} \| ^ {1 / n} = \inf _ {n > 0} \| x ^ {n} \| ^ {1 / n}
\]

称为 x 的谱半径。

对 A 的每个元素 \(x\)，我们有

\[
\begin{array}{r l} (3) & \varrho (x) \leqslant \| x \| \\ (4) & \varrho (x ^ {n}) = \varrho (x) ^ {n}, \text { for every integer } n \geqslant 1. \end{array}
\]

定义 3. --- A 的元素 x 称为拟幂零的，如果 \(\varrho(x)=0\)。

这等价于说，对任何 \(\lambda \in \mathbf{K}\)，诸数 \(\| (\lambda x)^n\|\) 对 \(n\geqslant 1\) 有界，或者等价地，对任何 \(\lambda \in \mathbf{K}\)，序列 \((\lambda x)^n\) 当 \(n\to +\infty\) 时趋于 0。

注 1. --- 设 A 是一个赋范代数。若元素 \(x \in \mathbf{A}\) 满足 \(\varrho(x) = \|x\|\)，则由 (3) 与 (4)，有 \(\|x^n\| = \|x\|^n\)（对一切 \(n \in \mathbf{N}\)）。

反之，假定有 \(\|x^{2}\|=\|x\|^{2}\)（对一切 \(x\in A\)）。则对每个整数 \(n\geqslant0\)，有等式 \(\|x^{2^{n}}\|=\|x\|^{2^{n}}\)，因而 \(\|x\|=\|x^{2^{n}}\|^{2^{-n}}\)；令 n 趋于 \(+\infty\)，就得到 \(\|x\|=\varrho(x)\)。

注 2. --- A 上的函数 \(x \mapsto \varrho(x)\)，作为连续函数 \(x \mapsto \|x^n\|^{1/n}\)（\(n \geqslant 1\)）的下包络，是上半连续的 (TG, IV, p.~31, 推论)，但一般说来它不连续。甚至可能发生这样的情形 (参见 I, p.~157, 习题 12)：A 的一个幂零元素序列收敛于一个不是拟幂零的元素。

\subsection{逆元}

设 A 是一个幺 Banach 代数，其单位元记作 1。回顾 (TG, IX, 命题 14, p.~40)，A 的可逆元所成的群 G 是 A 的一个开子集，A 的拓扑在 G 上诱导的拓扑与群结构相容，且拓扑群 G 是完备的。

命题 2. --- 设 A 是一个 Banach 代数，x 是 A 的一个元素。级数 \(\sum_{n=1}^{\infty}\lambda^{n}x^{n}\) 作为 \(\lambda\) 的幂级数来考虑，收敛半径为 \(\varrho(x)^{-1}\)。若 A 是幺的且 \(\varrho(x)<1\)，则 1-x 可逆，且其逆为 \(\sum_{n=0}^{\infty}x^{n}\)。

级数 \(\sum_{n=1}^{\infty}\lambda^{n}x^{n}\) 的收敛半径为

\[
(\limsup _ {n \to + \infty} \| x ^ {n} \| ^ {1 / n}) ^ {- 1} = \varrho (x) ^ {- 1}
\]

(参见 VAR, R1, p.~23, 3.1.4)。假定 A 有单位元。若 \(\varrho(x) < 1\)，则级数 \(\sum_{n=0}^{\infty} x^n\) 绝对收敛。由于

\[
(1 - x) \left(\sum_ {n = 0} ^ {k} x ^ {n}\right) = \left(\sum_ {n = 0} ^ {k} x ^ {n}\right) (1 - x) = 1 - x ^ {k + 1}
\]

对每个整数 \(k \geqslant 0\) 成立，元素 1 - x 可逆，且其逆等于 \(\sum_{n=0}^{\infty} x^{n}\)。

推论 1. --- 若 A 是一个幺 Banach 代数，则 A 的可逆元所成的群包含以 1 为中心、1 为半径的开球。

这是直接的，因为 \(\|x\| < 1\) 蕴涵 \(\varrho(x) < 1\)。

推论 2. --- 设 A 是一个 Banach 代数，I 是 A 的一个正则极大左（相应地，右）理想。则 I 是闭的。

设 ( \(\tilde{A}, e\) ) 是由 A 添加单位元而得到的幺 Banach 代数。存在 \(\tilde{A}\) 的一个极大左（相应地，右）理想 J，使得 \(J \cap A = I\) (A, VIII, p.~428, 命题 4)。于是 J 与以 e 为中心、1 为半径的开球不相交 (推论 1)，因而 \(\overline{J} \neq \widetilde{A}\)。由于 J 是极大理想，这就蕴涵 \(\overline{J} = J\)，从而 \(I = J \cap A = \overline{J} \cap A\) 在 A 中是闭的。

推论 3. --- Banach 代数的根是闭的。

事实上，根是诸正则极大左理想之交 (A, VIII, p.~430, 定义 3)。

命题 3. --- 设 A 是一个幺 Banach 代数。

\begin{enumerate}
\def\labelenumi{\alph{enumi})}
\item
  若 \(x \in A\) 有一个左逆（相应地，右逆）y，则每个元素 \(x' \in A\)（满足 \(\|x' - x\| < \|y\|^{-1}\)）都有一个左逆（相应地，右逆）。
\item
  A 的可逆（相应地，左可逆，相应地，右可逆）元素所成的集合在 A 中是开的。
\item
  设 \((x_{n})\) 是 A 的元素的一个序列，它们有左逆（相应地，右逆）\(y_{n}\)，使得 \((x_{n})\) 收敛于一个元素 \(x \in A\)，且序列 \((y_{n})\) 有界。则 \(x\) 是左可逆的（相应地，右可逆的）。
\end{enumerate}

只须处理左逆的情形；右逆的情形通过考虑反代数而得到。

设 \(x, y, x' \in \mathrm{A}\)，使得 \(yx = 1\) 且 \(\|x' - x\| < \|y\|^{-1}\)。我们有

\[
\left\| 1 - y x ^ {\prime} \right\| = \left\| y x - y x ^ {\prime} \right\| \leqslant \left\| y \right\| \cdot \left\| x - x ^ {\prime} \right\| <   1,
\]

因而 \(yx'\) 可逆：存在 \(z \in A\)，使得 \(z(yx') = 1\)。于是元素 \(x'\) 左可逆，其逆为 zy。这就证明了断言 a)，而断言 b) 立即得出。

  设 \((x_{n})\) 是 A 的元素的一个序列，它们有左逆 \(y_{n}\)，使得 \((x_{n})\) 收敛于一个元素 \(x \in A\)，且序列 \((y_{n})\) 有界。若 \(M \geqslant 1\) 是使 \(\|y_{n}\| \leqslant M\)（对一切 \(n \geqslant 1\)）的实数，则对充分大的 n，我们有 \(\|x_{n} - x\| < M^{-1} \leqslant \|y_{n}\|^{-1}\)，因而由 a)，x 有一个左逆。

定义 4. --- 设 A 是一个赋范代数，x 是 A 的一个元素。用 \(\gamma_x\) 与 \(\delta_x\) 表示从 A 到 A 的映射 \(y \mapsto xy\) 与 \(y \mapsto yx\)。我们说 \(x\) 是一个左（相应地，右）拓扑零因子，如果 \(\gamma_x\)（相应地，\(\delta_x\)）不是从 A 到 \(\gamma_x(A)\)（相应地，到 \(\delta_x(A)\)）上的一个同胚。

注. --- 由 TG, IX, p.~36, 推论 2，x 是一个左（相应地，右）拓扑零因子必须且只需存在 A 中的一个序列 \((z_{n})\)，使得 \(\|z_{n}\|=1\)，且使得 \(xz_{n}\)（相应地 \(z_{n}x\)）当 \(n\to+\infty\) 时趋于 0。

左（相应地，右）拓扑零因子是一个左（相应地，右）零因子。假定 A 非零且是幺的。左（相应地，右）拓扑零因子 x 不是左（相应地，右）可逆的。事实上，例如若 yx = 1 且 \(xz_{n}\) 趋于 0，则 \(z_{n} = y(xz_{n})\) 趋于 0，从而不可能对一切 n 有 \(\|z_{n}\| = 1\)。

命题 4. --- 设 A 是一个幺 Banach 代数。设 x 是 A 的一个非左可逆的元素。若存在 A 的左可逆元素的一个序列 \((x_{n})\)，它收敛于 x，则 x 是一个右拓扑零因子。

设 \(y_{n}\) 是 \(x_{n}\) 的一个左逆。由命题 3 (ii)，\(\|y_{n}\|\) 趋于 \(+\infty\)。令 \(z_{n} = \|y_{n}\|^{-1}y_{n}\)。我们有 \(\|z_{n}\| = 1\)，且 \(z_{n}x_{n} = \|y_{n}\|^{-1}\) 趋于 0，因而 \(z_{n}x = z_{n}x_{n} + z_{n}(x - x_{n})\) 趋于 0。于是借助定义 4 后面的注即可得出结论。

命题 5. --- 设 A 是一个幺 Banach 代数，B 是 A 的一个全子代数。则 \(\overline{\mathbf{B}}\) 是 A 的一个全子代数。

事实上，设 \(x\) 是 \(\overline{\mathrm{B}}\) 的一个在 A 中可逆的元素，并设 \((x_n)\) 是 B 的元素的一个序列，它趋于 \(x\)。则对充分大的 \(n\)，\(x_n\) 在 A 中可逆，且 \(x_n^{-1}\) 趋于 \(x^{-1}\)。由于子代数 B 是全的，我们有 \(x_n^{-1} \in \mathrm{B}\)，从而 \(x^{-1} \in \overline{\mathrm{B}}\)。

设 A 是一个幺 Banach 代数，\((y_{i})_{i\in\mathrm{I}}\) 是 A 的元素的一个族。设 B 是由诸元素 \(y_{i}\) 生成的 A 的全子代数。则 \(\overline{B}\) 是包含诸 \(y_{i}\) 的 A 的最小闭全子代数。它称为由诸元素 \(y_{i}\) 生成的闭全子代数。

\subsection{赋范代数中元素的谱}

在本小节中，我们假定 K = C。

定理 1. --- 设 A 是一个幺 Banach 代数，并设 \(x \in \mathrm{A}\)。

\begin{enumerate}
\def\labelenumi{\alph{enumi})}
\item
  集合 \(\mathrm{Sp}_{\mathrm{A}}(x)\) 是 C 的一个紧子集；
\item
  谱半径 \(\varrho(x)\) 是 \(\mathbf{C}\) 中以 0 为中心且包含 \(\mathrm{Sp}_{\mathrm{A}}(x)\) 的最小闭圆盘的半径；
\item
  预解式 \(\lambda \mapsto \mathrm{R}(x,\lambda) = (\lambda -x)^{-1}\)（\(x\) 的预解式）在 \(\mathbf{C} - \mathrm{Sp}_{\mathrm{A}}(x)\) 中全纯且在无穷远处为 0。而且，对每个整数 \(k\geqslant 0\)，我们有
下列公式：

\[
\frac {\partial^ {k}}{\partial \lambda^ {k}} \mathrm{R} (x, \lambda) = (- 1) ^ {k} k! \mathrm{R} (x, \lambda) ^ {k + 1};
\]

\item
  对每个复数 \(\lambda\)（满足 \(|\lambda| > 1 / \varrho(x)\)），有

\[
\mathrm{R} (x, \lambda) = \sum_ {n = 0} ^ {+ \infty} \lambda^ {- n - 1} x ^ {n}.
\]

\end{enumerate}

谱 \(\mathrm{Sp}_{\mathrm{A}}(x)\) 的补集是 A 的可逆元所成的群 G 在从 C 到 A 的连续映射 \(\lambda \mapsto x - \lambda\) 下的原像；由 I, p.~23 的命题 3, b)，谱 \(\mathrm{Sp}_{\mathrm{A}}(x)\) 是 C 的一个闭子集。此外，设 \(\lambda \in C\) 满足 \(|\lambda| > \varrho(x)\)。我们有 \(\lambda - x = \lambda(1 - \lambda^{-1}x)\)。由于 \(\varrho(\lambda^{-1}x) = |\lambda|^{-1}\varrho(x) < 1\)，元素 \(1 - \lambda^{-1}x\)，从而元素 \(\lambda - x\)，可逆，且

\[
\mathrm{R} (x, \lambda) = (\lambda - x) ^ {- 1} = \sum_ {n = 0} ^ {\infty} \lambda^ {- n - 1} x ^ {n}\tag{5}
\]

(I, p.~22, 命题 2)。特别地，\(\lambda \notin \mathrm{Sp}_{\mathrm{A}}(x)\)。这就证明 \(\mathrm{Sp}_{\mathrm{A}}(x)\) 包含在以 0 为中心、\(\varrho(x)\) 为半径的圆盘中。因而 \(\mathrm{Sp}_{\mathrm{A}}(x)\) 是紧的。公式 (5) 还表明，\(x\) 的预解式在闭圆盘 \(\Delta_{\varrho(x)}\)（以 0 为中心、\(\varrho(x)\) 为半径）的补集中有定义且全纯，并在无穷远处趋于 0。

设 \(\lambda_0\in \mathbf{C} - \mathrm{Sp}_{\mathrm{A}}(x)\)。令 \(y = \lambda_0 - x\)。设 \(\mu \in \mathbf{C}\) 满足 \(|\lambda_0 - \mu | < \| y^{-1}\|^{-1}\)。我们有

\[
\mu - x = y - (\lambda_ {0} - \mu) = y (1 - (\lambda_ {0} - \mu) y ^ {- 1}),
\]

因而 \(\mu - x\) 可逆，且其逆为

\[
(\mu - x) ^ {- 1} = y ^ {- 1} \sum_ {n = 0} ^ {\infty} (\lambda_ {0} - \mu) ^ {n} y ^ {- n}\tag{6}
\]

(由 I, p.~22 的命题 2)。于是 \(x\) 的预解式在以 \(\lambda_0\) 为中心、\(\| y^{-1} \|^{-1}\) 为半径的开圆盘中有定义且全纯。因而 \(x\) 的预解式是从 \(\mathbf{C} - \mathrm{Sp}_{\mathrm{A}}(x)\) 到 A 中的一个全纯映射。

I, p.~4 的公式 (1) 蕴涵 \(\frac{\partial}{\partial\lambda}\mathrm{R}(x,\lambda) = -\mathrm{R}(x,\lambda)^2\)，由此，对 \(k\) 用归纳法，

\[
\frac {\partial^ {k}}{\partial \lambda^ {k}} \mathrm{R} (x, \lambda) = (- 1) ^ {k} k! \mathrm{R} (x, \lambda) ^ {k + 1}.
\]

设 \(a > 0\) 是使 \(\mathrm{Sp}_{\mathrm{A}}(x)\) 包含在闭圆盘 \(\Delta_{a}\)（以 0 为中心、\(a\) 为半径）中的实数。函数 \(\lambda \mapsto (\lambda^{-1} - x)^{-1}\) 于是在 \(0 < |\lambda| < a^{-1}\) 上有定义且全纯，且当 \(\lambda\) 趋于 0 时趋于 0。在这个开圆盘（以 0 为中心、\(a^{-1}\) 为半径）上延拓此全纯函数的唯一连续函数是全纯的 (VAR, R1, 3.3.9)，因而定义它的级数 (5) 的收敛半径 \(\geqslant a^{-1}\) (VAR, R1, 3.2.9)。由 I, p.~22 的命题 2，我们因此有 \(a \geqslant \varrho(x)\)。

注. --- 1) 幺 Banach 代数中元素的谱可以是 C 的任意一个非空紧子集 F (参见 I, p. 17 的例 3；对 A = C(F; C) 与 f ∈ A（F 到 C 中的典范包含），有 \(\mathrm{Sp}_{\mathrm{A}}(f) = F\))。

\begin{enumerate}
\def\labelenumi{\arabic{enumi})}
\setcounter{enumi}{1}
\tightlist
\item
  设 A 是一个幺 Banach 代数，并设 \(x \in A\)。由 I, p.~24 的定理 1，\(\mathbf{C} - \mathrm{Sp}_{\mathrm{A}}(x)\) 是 C 的一个开子集，故是局部连通的。于是 \(\mathbf{C} - \mathrm{Sp}_{\mathrm{A}}(x)\) 的连通分支都是开的。由定理 1，这些连通分支之一包含由 \(\lambda \in C\)（满足 \(|\lambda| > \varrho(x)\)）所成的集合；所有其他的连通分支因此都是有界的。
\end{enumerate}

推论 1. --- 设 A 是一个非零的赋范幺代数。对每个 \(x \in \mathrm{A}\)，谱 \(\mathrm{Sp}_{\mathrm{A}}(x)\) 非空。

先假定 A 完备。若有 \(\mathrm{Sp}(x)=\varnothing\)，则 x 的预解式就会在 C 中全纯且在无穷远处为 0，从而恒等于零 (VAR, R., 3.3.6, p.~29)。由于 \(\mathrm{R}(x,\lambda)=(\lambda-x)^{-1}\) 可逆，就会得出 1=0，从而 \(A=\{0\}\)。

在一般情形，设 \(\widehat{\mathrm{A}}\) 是 A 的完备化代数；关系 \(\mathrm{Sp}_{\mathrm{A}}(x) = \varnothing\) 将蕴涵 \(\mathrm{Sp}_{\widehat{\mathrm{A}}} (x) = \varnothing\)，从而 \(\widehat{\mathrm{A}} = \{0\}\)，故 \(\mathrm{A} = \{0\}\)。

推论 2 (Gelfand–Mazur 定理). --- 设 A 是 C 上的一个赋范代数。若 A 是一个域，则 A = C · 1。

若 \(x \in A\)，则存在 \(\lambda \in \mathbf{C}\) 使 \(x - \lambda\) 不可逆 (推论 1)，因而 \(x - \lambda = 0\)，即 \(x \in \mathbf{C} \cdot 1\)。

推论 3. --- 设 A 是一个幺 Banach 代数，x 是 A 的一个可逆元素，满足 \(\|x\|=\|x^{-1}\|=1\)。则 \(\mathrm{Sp}(x)\subset\mathbf{U}\)。

设 \(\Delta\) 是 \(\mathbf{C}\) 中以 0 为中心、1 为半径的圆盘。由定理 1 b) 以及 \(\varrho(x) \leqslant \|x\|\) 这一事实，有 \(\operatorname{Sp}(x) \subset \Delta\)。同样，\(\operatorname{Sp}(x)^{-1} = \operatorname{Sp}(x^{-1}) \subset \Delta\)，由此得推论 (参见 I, p.~2, 注 4)。

推论 4. --- 设 E 是一个复 Banach 空间，\(\mathcal{L}(\mathrm{E})\) 是 E 的连续自同态所成的 Banach 代数，A 是 \(\mathcal{L}(\mathrm{E})\) 的一个非零子代数，使得 E 是一个单 A-伪模 (A, II, p.~176, 附录)。

\begin{enumerate}
\def\labelenumi{\alph{enumi})}
\item
  设 u 是 E 的一个自同态，不必连续，它与 A 交换。则 u 是一个位似。
\item
  设 \(u\) 是 \(\mathbf{E}\) 的一个自同态，不必连续。对每个整数 \(n \geqslant 1\) 及每个 \((\xi_1, \ldots, \xi_n) \in \mathrm{E}^n\)，存在 \(v \in A\)，使得

\[
(v (\xi_ {1}), \dots , v (\xi_ {n})) = (u (\xi_ {1}), \dots , u (\xi_ {n})).
\]

\end{enumerate}

我们来证明 \(a\) )。设 \(\widetilde{\mathrm{A}}\) 是由 A 添加单位元而得到的代数。由于 E 是一个单 A-伪模，它是一个单 \(\widetilde{\mathrm{A}}\)-模。

设 B 是 \(\tilde{\mathsf{A}}\) 在 C-向量空间 E 的自同态环中的交换子。代数 B 包含 1，且是 \(\tilde{\mathsf{A}}\)-模 E 的自同态代数。由于 E 是一个单 \(\tilde{\mathsf{A}}\)-模，Schur 引理 (A, VIII, p.~43, 推论) 表明 B 是一个域。

设 \(\xi_0\in \mathrm{E}\) 使 \(\mathrm{A}\xi_0\neq\{0\}\)。于是我们有 \(\mathrm{A}\xi_0 = E\)。对每个 \(u\in \mathrm{B}\)，用 \(\mathrm{A}_u\) 表示由 \(v\in \mathrm{A}\)（使 \(v(\xi_0) = u(\xi_0)\)）所成的集合。这个集合非空，因为 \(\mathrm{A}\xi_0 = E\)。我们令

\[
\| u \| _ {\mathrm{B}} = \inf _ {v \in \mathrm{A} _ {u}} \| v \|.
\]

映射 \(u \mapsto \|u\|_{\mathrm{B}}\) 是 B 上的一个半范数。

我们来证明这个映射是一个范数。设 u 是 B 的一个非零元素。对每个 \(v \in A_{u}\)，有 \(\|v\| \geqslant \|v(\xi_{0})\|/\|\xi_{0}\| = \|u(\xi_{0})\|/\|\xi_{0}\|\)，从而 \(\|u\|_{B} \geqslant \|u(\xi_{0})\|/\|\xi_{0}\|\)。于是只须证明 \(u(\xi_{0}) \neq 0\)。设 \(\xi_{1} \in E\) 使 \(u(\xi_{1}) \neq 0\)。由于 \(A\xi_{0} = E\)，存在 \(w \in A\) 使 \(\xi_{1} = w(\xi_{0})\)。于是 \(wu(\xi_{0}) = uw(\xi_{0}) = u(\xi_{1}) \neq 0\)，因而 \(u(\xi_{0}) \neq 0\)。

另一方面，设 \(u\) 与 \(u'\) 是 B 的元素。对每个 \(\varepsilon > 0\)，存在 \(v, v' \in A\)，使得 \(v(\xi_0) = u(\xi_0), v'(\xi_0) = u'(\xi_0)\)，且 \(\|v\| \leqslant \|u\|_{\mathrm{B}} + \varepsilon\)、\(\|v'\| \leqslant \|u'\|_{\mathrm{B}} + \varepsilon\)。于是我们有 \(vv'(\xi_0) = vu'(\xi_0) = u'v(\xi_0) = u'u(\xi_0)\)，从而

\[
\| u ^ {\prime} u \| _ {\mathrm{B}} \leqslant \| v ^ {\prime} v \| \leqslant \| v \| \| v ^ {\prime} \| \leqslant (\| u \| _ {\mathrm{B}} + \varepsilon) (\| u ^ {\prime} \| _ {\mathrm{B}} + \varepsilon),
\]

最后有 \(\|u'u\|_{B} \leqslant \|u\|_{B} \|u'\|_{B}\)。这表明 B 带上范数 \(u \mapsto \|u\|_{B}\)，是一个赋范代数。由于它是一个域，推论 2 蕴涵 \(B = C \cdot 1\)，这就是所要的结论。

我们证明 \(b\) )。由 \(a\) )，A 在 \(\mathrm{End}_{\mathbf{C}}(\mathrm{E})\) 中的交换子化为 E 的位似。于是其双交换子是 \(\mathrm{End}_{\mathbf{C}}(\mathrm{E})\)。断言 \(b\) ) 于是由 Jacobson 稠密定理 (A, VIII, p.~434 的定理 1) 得出。

推论 5. --- 设 A 是一个 Banach 代数，\(x \in A\)。

\begin{enumerate}
\def\labelenumi{\alph{enumi})}
\item
  \(\mathrm{Sp}^{\prime}(x)\) 是 C 的一个紧子集；
\item
  谱半径 \(\varrho(x)\) 是 C 中以 0 为中心且包含 \(\mathrm{Sp}'(x)\) 的最小闭圆盘的半径；
\item
  要使 \(x\) 是拟幂零的，必须且只需 \(\mathrm{Sp}'(x) = \{0\}\)。
\end{enumerate}

断言 \(a)\) 与 \(b)\) 由定理 1 通过考虑由 A 添加单位元而得到的 Banach 代数得出。断言 \(c)\) 由 \(b)\) 得出。

\subsection{相对于子代数的谱}

在本小节中，我们假定 K = C。

引理 3. --- 设 \(X_{1}\) 与 \(X_{2}\) 是 C 的紧子集。若 \(X_{2}\) 包含在 \(X_{1}\) 中，且 \(X_{1}\) 在 C 中的边界包含在 \(X_{2}\) 中，则 \(X_{1}\) 是 \(X_{2}\) 与 \(X_{2}\) 在 C 中的补集的某些有界连通分支的并。

设 U 是 \(\mathbf{C}-X_{2}\) 的一个连通分支。\(X_{1} \cap U\) 在开集 U 中的每个边界点也是 \(X_{1}\) 在 C 中的边界点，故由假设属于 \(X_{2}\)；由于 \(U \cap X_{2} = \varnothing\)，我们看出 \(X_{1} \cap U\) 在空间 U 中没有边界点。由于 U 连通，交 \(X_{1} \cap U\) 或为空或等于 U (TG, I, p.~82, 推论)，由此得引理。

命题 6. — 设 A 是一个幺 Banach 代数，B 是 A 的一个闭幺子代数。对每个 \(x \in B\)，有 \(\mathrm{Sp}_{\mathrm{B}}(x) \supset \mathrm{Sp}_{\mathrm{A}}(x)\)，且 \(\mathrm{Sp}_{\mathrm{A}}(x)\) 在 \(\mathbf{C}\) 中的边界包含 \(\mathrm{Sp}_{\mathrm{B}}(x)\) 在 \(\mathbf{C}\) 中的边界。特别地，若 \(\mathrm{Sp}_{\mathrm{B}}(x) \subset \mathbf{R}\)，则我们有 \(\mathrm{Sp}_{\mathrm{B}}(x) = \mathrm{Sp}_{\mathrm{A}}(x)\)。

我们有 \(\mathrm{Sp}_{\mathrm{B}}(x) \supset \mathrm{Sp}_{\mathrm{A}}(x)\) (I, p.~3 的注 6)。若 \(\lambda\) 是 \(\mathrm{Sp}_{\mathrm{B}}(x)\) 在 C 中的边界上的一个点，则存在一个序列 \((\lambda_{n})\)（由 \(\mathrm{Sp}_{\mathrm{B}}(x)\) 的外部点组成），它趋于 \(\lambda\)。于是 \(x - \lambda_{n}\) 在 B 中可逆，且趋于 \(x - \lambda\)，后者在 B 中不可逆；因而 \(x - \lambda\) 是 B 中的一个左或右拓扑零因子 (I, p.~24 的命题 4)，从而也是 A 中的。于是 \(\lambda \in \mathrm{Sp}_{\mathrm{A}}(x)\)。但由于 \(\mathrm{Sp}_{\mathrm{A}}(x) \subset \mathrm{Sp}_{\mathrm{B}}(x)\)，复数 \(\lambda \in \mathrm{Fr}_{\mathbf{C}}(\mathrm{Sp}_{\mathrm{B}}(x))\) 不可能是 \(\mathrm{Sp}_{\mathrm{A}}(x)\) 的内点，故属于其边界。

推论. --- 集合 \(\mathrm{Sp}_{\mathrm{B}}(x)\) 是 \(\mathrm{Sp}_{\mathrm{A}}(x)\) 与 \(\mathbf{C}-\mathrm{Sp}_{\mathrm{A}}(x)\) 的某些有界连通分支的并。

这由命题 6 与引理 3 得出。

这个推论将由 I, p.~46 的命题 13 与 I, p.~46 的命题 14 加以补充。

\section{交换 Banach 代数}

在本节中，基域是 C。

\subsection{交换 Banach 代数的特征标}

定理 1. --- 设 A 是一个 Banach 代数，并设 \(\chi\colon\mathbf{A}\to\mathbf{C}\) 是一个代数态射 (参见 I, p.~9)。则 \(\chi\) 连续，且范数至多为 1。若 A 是幺的且 \(\chi\) 是一个保单位元态射，则 \(\chi\) 的范数为 1。

我们来证明 \(|\chi(x)| \leqslant \|x\|\)（对一切 \(x \in \mathbf{A}\)）。如有必要，把 A 换成由 \(x\) 生成的 Banach 代数，我们就可以假定 A 是交换的；这时 \(\chi \in \mathbf{X}'(\mathbf{A})\)。对每个 \(x \in \mathbf{A}\)，有 \(\chi(x) \in \mathrm{Sp}_{\mathrm{A}}'(x)\) (I, p.~9, 第 7 小节)，因而 \(|\chi(x)| \leqslant \varrho(x) \leqslant \|x\|\) (I, p.~28, 推论 5)，由此得第一个断言。

此外，若 A 是幺的，则等式 \(\chi(1)=1\) 蕴涵 \(\|\chi\|\geqslant1\)，由此得所要的等式。

注. --- 设 A 是一个交换 Banach 代数。由这个定理得出，\(X'(A)\) 上的逐点收敛拓扑与 \(X'(A)\) 上的诱导拓扑重合，后者由弱拓扑 \(\sigma(A', A)\)（A 的对偶 \(A'\) 的弱拓扑）诱导而来。

推论. --- 设 A 是一个交换 Banach 代数。空间 X'(A) 是紧的。空间 X(A) 是局部紧的，且若 A 有单位元，则它是紧的。

设 A' 是 Banach 空间 A 的对偶。它的单位球 \(A_{1}^{\prime}\) 是弱紧的 (EVT, III, p.~17, 推论 2)。由定理 1，我们有 \(\mathsf{X}^{\prime}(\mathsf{A}) \subset \mathsf{A}_{1}^{\prime}\)。此外，\(\mathsf{X}^{\prime}(\mathsf{A})\) 在 A' 中对于弱拓扑是闭的，因为它是下列弱闭集的交

\[
\mathrm{X} _ {x, y} = \left\{f \in \mathrm{A} ^ {\prime} \mid f (x y) - f (x) f (y) = 0 \right\} \quad (\text { for } x, y \in \mathrm{A}).
\]

由此得出 \(X'(A)\) 是紧的，且 \(X(A) = X'(A) - \{0\}\) 是局部紧的。

若 A 有单位元 1，则 \(\mathsf{X}(\mathsf{A})\) 是由 \(\chi \in \mathsf{X}'(\mathsf{A})\)（满足 \(\chi(1) = 1\)）所成的集合，因此是 \(\mathsf{X}'(\mathsf{A})\) 的一个闭子集，从而是紧的。

每个紧空间都同胚于某个适当的交换幺 Banach 代数 A 的 X(A) (参见 I, p.~32, 推论 2)。

定理 2. --- 设 A 是一个交换 Banach 代数。映射 \(\chi \mapsto \operatorname{Ker}(\chi)\) 是从 \(\mathsf{X}(\mathsf{A})\) 到 A 的正则极大理想所成的集合 \(\mathrm{J}(\mathrm{A})\) 上的一个双射。

设 I 是 A 的一个正则极大理想。它是闭的 (I, p.~22, 推论 2)，因而 A/I 是一个 Banach 代数。由于它是一个域 (A, VIII, p.~426, 命题 2)，Gelfand--Mazur 定理 (I, p.~26, 推论 2) 蕴涵 A/I 在 C 上维数为 1。于是 I 在 A 中的余维数为 1。定理于是由 I, p.~10 的引理 3 得出。

可能发生这样的情形：一个非零的、无单位元的交换 Banach 代数 A 没有极大理想；这时 \(X'(A)\) 缩为 \(\{0\}\) (参见 I, p.~186, 习题 31)。

这个定理表明，I, p.~11 中的集合 \(J(A)\) 与 \(\hat{A}\)（由于 A 是交换的，这两者视为同一）也可以与集合 \(X(A)\) 视为同一。因此有理由在 \(X(A)\) 上考虑弱拓扑和 Jacobson 拓扑，后者更粗 (I, p.~15, 命题 5)。Jacobson 拓扑的闭集 \(V(M)\)（对 \(M \subset A\)）是下列集合

\[
\{\chi \in \mathrm{X} (\mathrm{A}) \mid \chi (\mathrm{M}) = 0 \}
\]

（\(\mathsf{X}(\mathsf{A})\) 的子集），它们有时仍记作 \(\mathsf{V}(\mathsf{M})\)。类似地，对任一子集 \(\mathsf{M}\)（\(\mathsf{X}(\mathsf{A})\) 的子集），我们用 \(\Upsilon(\mathsf{M})\) 表示诸 \(\chi \in \mathsf{M}\) 的核之交这个理想；它等于 I, p.~13 中所定义的集合 \(\Upsilon(\mathsf{M}')\)，即对子集 \(\mathsf{M}' \subset \mathsf{J}(\mathsf{A})\)（当把 \(\mathsf{J}(\mathsf{A})\) 与 \(\mathsf{X}(\mathsf{A})\) 视为同一时与 M 对应的子集）所定义的那个集合。

当在 X(A) 中使用一个拓扑概念而没有指明所用的是哪个拓扑时，所说的始终是弱拓扑。

\subsection{局部紧空间上在无穷远处为 0 的连续函数}

在本小节中，X 是一个局部紧空间。我们用 \(\mathcal{C}_{0}(X)\) 表示 X 上在无穷远处为 0 的复值连续函数所成的交换 Banach 代数，带上范数 \(\|f\|=\sup_{x\in X}|f(x)|\) (I, p.~17 的例 3)。

命题 1. --- 对 X 的每个闭子集 \(\Phi\)，用 \(I_{\Phi}\) 表示由 \(f \in \mathcal{C}_0(X)\)（在 \(\Phi\) 上为 0）所成的集合。则 \(\Phi \mapsto I_{\Phi}\) 是从 X 的闭子集所成的集合到 \(\mathcal{C}_0(X)\) 的闭理想所成的集合上的一个双射。

集合 \(\mathrm{I}_{\Phi}\) 是 \(\mathcal{C}_0(\mathrm{X})\) 的一个闭理想。

设 \(\Phi \neq \Phi'\) 是 X 的两个闭子集。如有必要，交换 \(\Phi\) 与 \(\Phi'\)，我们就可以假定存在 \(x \in \Phi'\) 使 \(x \notin \Phi\)；于是存在函数 \(f \in \mathcal{C}_0(X)\)，它在 \(\Phi\) 上为 0 且在 \(x\) 处非零 (TG, IX, p.~43, 命题 1)。我们有 \(f \in I_{\Phi}\) 且 \(f \notin I_{\Phi'}\)，因而映射 \(\Phi \mapsto I_{\Phi}\) 是单射。

设 I 是 \(\mathcal{C}_{0}(\mathrm{X})\) 的一个闭理想。设 \(\Phi\) 是由 \(x \in X\)（使 \(f(x) = 0\) 对一切 \(f \in I\) 成立）所成的集合；它是 X 的一个闭子集，且我们有 \(I \subset I_{\Phi}\)。我们来证明 \(I_{\Phi} \subset I\)，这将蕴涵 \(I = I_{\Phi}\) 并完成命题的证明。

设 \(f \in \mathrm{I}_{\Phi}\)。对每个实数 \(\varepsilon > 0\)，用 \(\mathrm{C}_{\varepsilon}\) 表示由 \(x \in \mathrm{X}\)（满足 \(|f(x)| \geqslant \varepsilon\)）所成的集合。由于 \(f\) 在无穷远处趋于 0，集合 \(\mathrm{C}_{\varepsilon}\) 是紧的。设 \(x \in \mathrm{C}_{\varepsilon}\)；由于 \(f(x) \neq 0\) 且 \(f \in \mathrm{I}_{\Phi}\)，有 \(x \notin \Phi\)；于是由 \(\Phi\) 的定义，存在函数 \(\varphi_x \in \mathrm{I}\) 使 \(|\varphi_x(x)| > 1\)，从而使得 \(|\varphi_x(y)| > 1\) 对一切 \(y\)（属于 \(\mathrm{V}_x\)，即 \(x\) 的一个邻域）成立。开集 \(\mathrm{V}_x \cap \mathrm{C}_{\varepsilon}\) 覆盖 \(\mathrm{C}_{\varepsilon}\)。由于集合 \(\mathrm{C}_{\varepsilon}\) 是紧的，存在有限子集 \(\mathrm{T}_{\varepsilon} \subset \mathrm{X}\)，使得

\[
\mathrm{C} _ {\varepsilon} \subset \bigcup_ {x \in \mathrm{T} _ {\varepsilon}} \mathrm{V} _ {x}.
\]

于是，元素

\[
g _ {\varepsilon} = \frac {1}{\varepsilon} \sum_ {x \in \mathrm{T} _ {\varepsilon}} \varphi_ {x} \overline {{\varphi_ {x}}} \geqslant 0
\]

（\(\mathcal{C}_0(\mathrm{X})\) 中的）属于 I，且我们有 \(g_{\varepsilon} \geqslant \varepsilon^{-1}\)（在 \(\mathrm{C}_{\varepsilon}\) 上）。函数

\[
f _ {\varepsilon} = \frac {f g _ {\varepsilon}}{1 + g _ {\varepsilon}}
\]

属于 I。对 \(x \notin \mathrm{C}_{\varepsilon}\)，有

\[
\left| f (x) - f _ {\varepsilon} (x) \right| \leqslant 2 \varepsilon ,
\]

而对 \(x\in \mathrm{C}_{\varepsilon}\)，有

\[
| f (x) - f _ {\varepsilon} (x) | = \frac {| f (x) |}{1 + g _ {\varepsilon} (x)} \leqslant \varepsilon | f (x) |.
\]

于是 \(f_{\varepsilon}\) 当 \(\varepsilon\) 趋于 0 时在 X 上一致收敛于 f。因此我们有 \(f \in \bar{I}\)，而由于 I 是闭的，就有 \(f \in I\)。

推论 1. --- 对每个 \(x \in X\)，用 \(I_x\) 表示由 \(f \in \mathcal{C}_0(X)\)（在 \(x\) 处为 0）所成的集合。则 \(x \mapsto I_x\) 是从 \(X\) 到 \(\mathcal{C}_0(X)\) 的极大闭理想所成的集合上的一个双射。这些理想都是正则的。

这立即由命题 1 得出。

用 X' 表示 X 的 Alexandroff 紧化，即由 X 添加一个无穷远点 \(\omega_{X}\) 而得到的紧空间 (TG, I, pp. 67–68)。代数 \(\mathcal{C}_{0}(X)\) 与 X' 上在 \(\omega_{X}\) 处为 0 的复值连续函数所成的 Banach 代数视为同一。

对每个 \(x \in X'\)，我们用 \(ev_{x}\) 表示 \(\mathcal{C}_{0}(X)\) 的由 \(\mathrm{ev}_{x}(f) = f(x)\)（对一切 \(f \in \mathcal{C}_{0}(X)\)）所定义的特征标。

推论 2. --- 映射 \(x \mapsto \mathrm{ev}_x\) 是从 \(\mathbf{X}'\) 到 \(\mathsf{X}'(\mathcal{C}_0(\mathbf{X}))\) 上的一个同胚，而它在 \(\mathbf{X}\) 上的限制是从 \(\mathbf{X}\) 到 \(\mathsf{X}(\mathcal{C}_0(\mathbf{X}))\) 上的一个同胚。此外，弱拓扑与 Jacobson 拓扑在 \(\mathsf{X}(\mathcal{C}_0(\mathbf{X}))\) 上重合。

映射 ev：\(x \mapsto ev_{x}\)（从 \(X'\) 到 \(\mathsf{X}'(\mathcal{C}_{0}(\mathbf{X}))\) 中的映射）是单射。由 I, p.~30 的推论 1 和定理 2，它是满射。它是连续的，因为对每个函数 \(f \in \mathcal{C}_{0}(\mathbf{X})\) 及 R 的每个开集 U，我们有

\[
\mathrm{ev} ^ {- 1} (\{\chi \in \mathsf {X} ^ {\prime} (\mathcal {C} _ {0} (\mathsf {X})) \mid \chi (f) \in \mathrm{U} \}) = f ^ {- 1} (\mathrm{U})
\]

它在 X 中是开的。由于 \(X'\) 是紧的，映射 ev 因此是一个同胚。于是 ev 在 X 上的限制是一个到 \(\mathsf{X}(\mathcal{C}_{0}(\mathsf{X}))\) 上的同胚。

若 F 是 \(\mathsf{X}(\mathcal{C}_{0}(\mathbf{X}))\) 的一个弱闭子集，则它经由同胚 ev 对应于 X 的一个闭子集 \(\Phi\)；确切地说，由命题 1，我们有 \(\mathrm{F} = \{\chi \in \mathsf{X}(\mathcal{C}_{0}(\mathbf{X})) \mid \mathrm{I}_{\Phi} \subset \mathrm{Ker} \, \chi\}\)，它对于 Jacobson 拓扑是闭的。

推论 3. --- 假定 X 是紧的。则映射 \(x \mapsto \mathrm{ev}_x\) 是从 X 到 \(\mathsf{X}(\mathcal{C}(\mathbb{X}))\) 上的一个同胚。弱拓扑与 Jacobson 拓扑在 \(\mathsf{X}(\mathcal{C}(\mathbb{X}))\) 上重合。
